%% ===================================================================
%% sections_leximin.tex --- plain LaTeX fragments (no preamble).
%% Assumes amsmath/amssymb/amsthm, booktabs, enumitem, natbib;
%% theorem-like envs {theorem,lemma,proposition,corollary,definition,
%% example,remark}; macros \R, \core.  Player i = index i (1-indexed).
%% Companion machine-readable certificate: supplement_leximin.json.
%% ===================================================================

\section{The leximin-above-floors core selection}\label{sec:lx-def}

Let $(N,v)$ be a balanced TU game, $M=v(N)$, and floors $m_i:=v(\{i\})$. For
$x\in\R^N$ write the \emph{gains} $g_i(x)=x_i-m_i$. The
\emph{leximin-above-floors} selection is
\[
  \phi(v):=\arg\operatorname*{leximin}_{x\in\core(v)}\bigl(g_1(x),\dots,g_n(x)\bigr),
\]
the unique core allocation whose ascending-sorted gain vector is
lexicographically largest. It is single-valued (leximin over a polytope) and is
computed by the Maschler-type scheme over player coordinates: repeatedly raise a
common lower level $s$ of the free players' gains as far as the core allows, fix
(``pin'') every free player whose gain cannot exceed $m_i+s$ on the level-optimal
face, and recurse on the rest. On \emph{exact} games $m_i=\min_{\core}x_i$, so
$g_i\ge 0$ and $\phi$ egalitarianly distributes the surplus $x_i-v(\{i\})$ inside
the core; on convex games it is the constrained egalitarian solution of
\citet{DuttaRay89}. Implementation (exact rational + float) in
\texttt{src/leximin.py}, validated against those two coincidences.

We ask whether $\phi$ is \emph{coalitionally monotonic}: if $v'=v+d\,\mathbf 1_S$
for a single coalition $S$ and $d>0$, is $\phi_i(v')\ge\phi_i(v)$ for every
$i\in S$? A positive answer on exact games would exhibit an explicit coalitionally
monotonic core selection on that class. The reduction in
Section~\ref{sec:lx-reduction} addresses only proper nonsingleton coalitions.
Singleton bumps change the floors defining $\phi$, and a grand-coalition bump
changes the efficiency hyperplane.  The singleton case is resolved next.  The
grand-coalition case has a negative answer even on a simple exact five-player
game (Theorem~\ref{thm:lx-grand-fails}).


\section{Singleton bumps are monotone}\label{sec:lx-singleton}

We first isolate an elementary property of the leximin order.  Write
$a\succ_{\rm lx}b$ when the ascending rearrangement of $a$ is
lexicographically larger than that of $b$.

\begin{lemma}[one-coordinate translation]\label{lem:lx-translation}
If $a\succ_{\rm lx}b$, $a_i<b_i$, and $d>0$, then
$a+d e_i\succ_{\rm lx}b+d e_i$.
\end{lemma}

\begin{proof}
For a finite vector $z$, put $L_\lambda(z)=\sum_j e^{-\lambda z_j}$.
The first unequal coordinate of the ascending rearrangements dominates this
sum as $\lambda\to\infty$; hence
$a\succ_{\rm lx}b$ if and only if
$L_\lambda(a)<L_\lambda(b)$ for all sufficiently large $\lambda$.
Moreover,
\begin{align*}
 L_\lambda(b+d e_i)-L_\lambda(a+d e_i)
 ={}&L_\lambda(b)-L_\lambda(a)\\
 &+(1-e^{-\lambda d})(e^{-\lambda a_i}-e^{-\lambda b_i}).
\end{align*}
Both terms on the right are positive for all sufficiently large $\lambda$,
because $a_i<b_i$.  The same characterization gives the claim.
\end{proof}

\begin{theorem}[singleton coalitional monotonicity]
\label{thm:lx-singleton}
Let $v$ and $v'=v+d\,\mathbf 1_{\{i\}}$ be balanced games, with $d>0$.
Then
\[
 \phi_i(v')\geq\phi_i(v).
\]
Thus singleton bumps are monotone even without exactness.
\end{theorem}

\begin{proof}
Let $m_j=v(\{j\})$ and express the old core in gain coordinates
\[
 G=\{(x_j-m_j)_{j\in N}:x\in\core(v)\}.
\]
If the singleton floor of player $i$ is raised by $d$, the new feasible gain
set is
\[
 H=\{g-d e_i:g\in G,\ g_i\geq d\}.
\]
Let $z$ be the leximin point of $G$ and $h$ the leximin point of $H$.
The desired inequality is $h_i+d\geq z_i$.  If $z_i<d$, this is immediate
from $h_i\geq0$.  Suppose $z_i\geq d$.  Then $b=z-d e_i$ belongs to $H$, so
$h$ is either $b$ or $h\succ_{\rm lx}b$.  If, contrary to the claim,
$h_i+d<z_i$, then $h_i<b_i$ and the equality case is impossible.
Lemma~\ref{lem:lx-translation} gives
\[
 h+d e_i\succ_{\rm lx}b+d e_i=z.
\]
But $h+d e_i\in G$, contradicting the definition of $z$ as the leximin point
of $G$.
\end{proof}


\section{Grand-coalition bumps: an exact path but no resource monotonicity}
\label{sec:lx-grand}

The grand-coalition case also has an exact path, but its geometry differs from
the nested-core path of Section~\ref{sec:lx-reduction}.

\begin{lemma}[grand-bump exact path]\label{lem:lx-grand-path}
Let $v$ be exact and, for $r\geq0$, let $v^r$ agree with $v$ on every proper
coalition and satisfy $v^r(N)=v(N)+r$.  Then $v^r$ is exact for every $r\geq0$.
Moreover,
\[
 \core(v)+r\Delta_N\subseteq\core(v^r),
 \qquad
 \Delta_N=\{a\in\R^N_+:a(N)=1\}.
\]
\end{lemma}

\begin{proof}
If $x\in\core(v)$ and $a\in r\Delta_N$, then $x+a$ has total
$v(N)+r$ and weakly raises every coalition sum, proving the inclusion and
balancedness.  Fix a proper coalition $T$.  Exactness of $v$ supplies
$x^T\in\core(v)$ with $x^T(T)=v(T)$.  Choose $j\notin T$.  Then
$x^T+r e_j\in\core(v^r)$ by the inclusion and remains tight at $T$.
Thus
\[
 \min_{x\in\core(v^r)}x(T)=v(T)=v^r(T).
\]
The grand coalition is tight throughout its core by efficiency, so $v^r$ is
exact at every coalition.
\end{proof}

The exact path does not rescue labeled resource monotonicity.

\begin{theorem}[grand resource monotonicity fails on exact games]
\label{thm:lx-grand-fails}
There is an exact five-player game $v$ such that, when only $v(N)$ is raised,
one player's leximin-above-floors payoff strictly falls.  Indeed, let
\[
 p=(0,0,7,0,7),\qquad q=(2,10,0,2,0),
\]
put $v(N)=14$, and for every proper $T$ put
$v(T)=\min\{p(T),q(T)\}$.  If $v^r(N)=14+r$ and $v^r(T)=v(T)$ for
$T\subsetneq N$, then for $0\leq\varepsilon\leq3/2$,
\[
 \phi(v^{2+\varepsilon})
 =\left(3+\frac{\varepsilon}{3},
          4-\frac{\varepsilon}{3},
          3+\frac{\varepsilon}{3},
          3+\frac{\varepsilon}{3},
          3+\frac{\varepsilon}{3}\right).
\]
Thus player $2$ falls at rate $1/3$ while the grand worth rises.  For example,
raising the grand worth from $16$ to $129/8$ changes that player's payoff from
$4$ to $95/24$.
\end{theorem}

\begin{proof}
Both $p$ and $q$ have total $14$ and belong to $\core(v)$.  At every proper
coalition one of them attains $v(T)$, so $v$ is exact; every $v^r$ is exact by
Lemma~\ref{lem:lx-grand-path}.  All singleton worths are zero.

Write $s=3+\varepsilon/3$ and
$z=(s,7-s,s,s,s)$.  Direct substitution in the lower-envelope definition of
$v$ shows that $z\in\core(v^{2+\varepsilon})$ throughout the stated
interval.  (At $\varepsilon=0$ every core slack is nonnegative; the only
coalition sum with negative slope is the singleton of player $2$, whose slack
remains positive.  Every other slope is nonnegative.)

It remains to identify the leximin point.  The two pair constraints
\[
 x_2+x_3\geq7,\qquad x_2+x_5\geq7
\]
belong to the core.  If every coordinate of a core point is at least $t$, then
efficiency and the first constraint give
$16+\varepsilon=x(N)\geq7+3t$, hence $t\leq s$.
So $z$ attains the largest possible minimum coordinate.  On that
level-optimal face put $d_i=x_i-s\geq0$ and
$D=1-2\varepsilon/3\geq0$.  Efficiency gives $\sum_i d_i=D$, while the two
displayed constraints give $d_2+d_3\geq D$ and $d_2+d_5\geq D$.
Nonnegativity forces $d_2=D$ and every other $d_i=0$.  Hence the
level-optimal face is the singleton $\{z\}$, proving the formula and the
strict drop.
\end{proof}

The counterexample pinpoints the failure mechanism: four players form the
lowest leximin stratum, and additional total resource raises them jointly only
by taking payoff from a higher-stratum player pinned to two cross-stratum pair
constraints.  The feasible comparison points supplied by
$\core(v)+r\Delta_N$ therefore do not control labeled coordinates.


\section{Exactness is necessary}\label{sec:lx-necessity}

Coalitional monotonicity of $\phi$ fails on balanced games in general---already at
$n=4$, and even at $n=5$ where every floor is core-attained. So exactness cannot be
weakened to balancedness.

\begin{theorem}\label{thm:lx-necessity}
There is a balanced (non-exact) five-player game $v$, a coalition $S$ with $|S|=3$,
and $d>0$ with $v+d\,\mathbf 1_S$ balanced, such that some $i\in S$ has
$\phi_i(v+d\,\mathbf 1_S)<\phi_i(v)$. The same holds for a four-player game.
\end{theorem}

\begin{example}[$n=5$; floors attained]\label{ex:lx-n5}
Let $N=\{1,\dots,5\}$ with $v(N)=7$ and (all other coalitions as listed in the
supplement) in particular
\[
  v(\{1,2,3\})=3,\quad v(\{2,3\})=3,\quad v(\{1,3,5\})=v(\{3,5\})=4,\quad
  v(\{1,5\})=-1 .
\]
Floors are $m=(0,1,1,1,2)$, and $m_i=\min_{\core}x_i$ for \emph{every} $i$ (all
floors attained); the game is balanced but not exact (e.g.\
$\min_{\core}x(\{1,5\})=1>-1$). Then
\[
  \phi(v)=\bigl(\tfrac12,\tfrac32,\tfrac32,1,\tfrac52\bigr),\qquad
  \phi(v+\mathbf 1_{\{1,2,3\}})=(1,1,2,1,2),
\]
both in their respective cores, and player $2\in\{1,2,3\}$ \emph{drops} from
$\tfrac32$ to $1$ when $v(\{1,2,3\})$ is raised from $3$ to $4$. Thus exactness at
the singletons (floor attainment) does not suffice: it is exactness at higher
coalitions that the positive theorem needs.
\end{example}

\begin{example}[$n=4$; floor unattained]\label{ex:lx-n4}
Let $N=\{1,2,3,4\}$, $v(N)=20$, with $v(\{1\})=-3$, $v(\{1,2\})=11$,
$v(\{1,3\})=7$, $v(\{2,3\})=4$, $v(\{1,2,3\})=13$, $v(\{2,4\})=10$,
$v(\{1,3,4\})=9$, $v(\{2,3,4\})=13$, and the remaining values in the supplement.
Here $v(\{1\})=-3<1=\min_{\core}x_1$. Then $\phi(v)=(\tfrac52,\tfrac{17}2,
\tfrac92,\tfrac92)$ and $\phi(v+3\,\mathbf 1_{\{1,2,3\}})=(2,9,5,4)$: player
$1\in\{1,2,3\}$ drops from $\tfrac52$ to $2$.
\end{example}

Both examples are certified in exact rational arithmetic ($\phi$ in the core on
both sides, the drop exact); see the supplement and \texttt{src/hunt\_leximin\_nonexact.py}.


\section{Reduction for proper nonsingleton bumps on exact games}\label{sec:lx-reduction}

Fix an exact $v$, a proper coalition $S$ with $|S|\ge 2$, and $d>0$ with
$v'=v+d\,\mathbf 1_S$ exact. Since $|S|\ge2$ the floors are unchanged, and
\[
  \core(v')=\core(v)\cap\{x:x(S)\ge v(S)+d\}\subseteq\core(v).
\]
For $\tau\in[v(S),v(S)+d]$ let $v_\tau$ agree with $v$ except $v_\tau(S)=\tau$, so
$\core(v_\tau)=\core(v)\cap\{x(S)\ge\tau\}$ interpolates $\core(v)$ and $\core(v')$.

\begin{lemma}[exact interval]\label{lem:lx-interval}
Every $v_\tau$, $\tau\in[v(S),v(S)+d]$, is exact.
\end{lemma}

\begin{proof}
$\core(v_\tau)\supseteq\core(v')\neq\emptyset$, so $v_\tau$ is balanced.
Exactness needs $\min_{\core(v_\tau)}x(T)=v_\tau(T)$ for all $T$.
For $T=S$, exactness of $v$ supplies $x^0\in\core(v)$ with
$x^0(S)=v(S)$, while any $x^1\in\core(v')$ satisfies
$x^1(S)\ge v(S)+d$. The core of $v$ is convex, so the segment
$[x^0,x^1]$ contains a point $x^\tau\in\core(v)$ with $x^\tau(S)=\tau$.
This point belongs to $\core(v_\tau)$, proving
$\min_{\core(v_\tau)}x(S)=\tau$.
For $T\neq S$: put $f_T(\tau)=\min_{\core(v_\tau)}x(T)$. As $\tau$ grows
$\core(v_\tau)$ shrinks, so $f_T$ is nondecreasing; and $\core(v_\tau)\subseteq
\core(v)$ gives $f_T(\tau)\ge\min_{\core(v)}x(T)=v(T)$. At the endpoints
$f_T=v(T)$ because $v$ and $v'$ are exact. A nondecreasing function equal to
$v(T)$ at both endpoints and $\ge v(T)$ between is $\equiv v(T)$. Hence
$f_T(\tau)=v(T)=v_\tau(T)$.
\end{proof}

The entire exact-bump range can be read from exposed core faces. For
$T\subseteq N$, write
$F_T=\{x\in\core(v):x(T)=v(T)\}$ (with $F_N=\core(v)$), and define
\[
 d^*_v(S)=\min_{T\neq S}
 \left(\max_{x\in F_T}x(S)-v(S)\right).
\]

\begin{proposition}[exact bump capacity]\label{prop:lx-capacity}
For an exact game $v$ and a proper coalition $S$, the bumped game
$v+d\mathbf 1_S$ is exact if and only if $0\leq d\leq d^*_v(S)$.
If $\sigma_S=\phi(v)(S)-v(S)$, then every exact bump at $S$ is
non-binding if and only if $d^*_v(S)\leq\sigma_S$; a binding exact bump
exists if and only if $d^*_v(S)>\sigma_S$.
\end{proposition}

\begin{proof}
The bumped core is
$K_d=\core(v)\cap\{x:x(S)\geq v(S)+d\}$. For $T\neq S$, exactness at
$T$ is equivalent to $F_T\cap K_d\neq\emptyset$, which is equivalent to
\[
 d\leq \max_{x\in F_T}x(S)-v(S).
\]
Taking the minimum over $T\neq S$ gives all the higher-coalition exactness
conditions; the term $T=N$ also gives nonemptiness of $K_d$. Exactness at
$S$ follows by interpolating between a point of $F_S$ and a point of $K_d$,
as in Lemma~\ref{lem:lx-interval}. This proves the first statement. The
last two statements compare the exact capacity with the largest bump still
satisfied by the old selection, namely $\sigma_S$.
\end{proof}

The capacity implication can be proved for every star-failure dual with at
most two residual levels.  This includes every two-player coalition and also
explains the same-stratum pair certificates found in the computation.

\begin{proposition}[two-level dual certificate]
\label{prop:lx-two-level-capacity}
Let $v$ be exact, $S$ proper, $\phi=\phi(v)$, and
$\sigma_S=\phi(S)-v(S)$.  Suppose $(\star_S)$ fails and its dual maximal
face can be represented by weights taking at most two values on $S$: after
normalization, for some $\emptyset\neq C\subsetneq S$ and $R>1$,
\[
 \phi\in\arg\max_{x\in\core(v)}
 \left(x(C)+R x(S\setminus C)\right).
\]
Then
\[
 \max_{x\in F_C}x(S)\leq\phi(S)
 \quad\text{and hence}\quad
 d_v^*(S)\leq\sigma_S.
\]
The same conclusion holds in the one-level case where $\phi$ maximizes
$x(S)$.
\end{proposition}

\begin{proof}
Put $H=S\setminus C$ and $h=\phi(C)-v(C)\geq0$.  Exactness makes
$F_C=\{x\in\core(v):x(C)=v(C)\}$ nonempty.  For $x\in F_C$, maximality
of $\phi$ gives
\[
 (x(C)-\phi(C))+R(x(H)-\phi(H))\leq0.
\]
The first parenthesis is $-h$, so
$x(H)-\phi(H)\leq h/R$.  Therefore
\[
 x(S)-\phi(S)
 =-h+x(H)-\phi(H)
 \leq-h\left(1-\frac1R\right)\leq0.
\]
Thus $F_C$ certifies the required bound:
\[
 d_v^*(S)\leq\phi(S)-v(S)=\sigma_S.
\]
In the one-level case, $\phi(S)=\max_{\core(v)}x(S)$, so the conclusion is
immediate.
\end{proof}

\begin{corollary}[the capacity implication for pairs]
\label{prop:lx-pair-capacity}
If $v$ is exact and $S$ is a proper two-player coalition, then
\[
 d_v^*(S)>\sigma_S\quad\Longrightarrow\quad(\star_S).
\]
\end{corollary}

\begin{proof}
When $(\star_S)$ fails, its positive dual weights on the two members of $S$
have either one value or two.  Normalize the smaller value to one and apply
Proposition~\ref{prop:lx-two-level-capacity}.
\end{proof}

There is a second exact certificate, independent of the number of residual
levels.  It compares the slack of a coalition carrying the whole complement
of $S$ with the slack of its trace inside $S$.

\begin{proposition}[off-side slack-domination certificate]
\label{prop:lx-carrier-capacity}
Let $v$ be exact, $S$ proper, $\phi=\phi(v)$, and $P=N\setminus S$.
Suppose $C$ is a nonempty proper subset of $S$, put $T=P\cup C$, and assume
\[
 \phi(T)-v(T)\leq\phi(C)-v(C).
\]
Then
\[
 \max_{x\in F_C}x(S)\leq\phi(S),
 \qquad\text{so}\qquad
 d_v^*(S)\leq\phi(S)-v(S).
\]
In particular, the conclusion holds whenever $T$ is $\phi$-tight.
\end{proposition}

\begin{proof}
Because $S\cup T=N$ and $S\cap T=C$, modularity gives
$x(S)+x(T)=x(N)+x(C)$ for every allocation $x$.  If $x\in F_C$, then
core feasibility at $T$ and efficiency imply
\[
 x(S)=v(N)+v(C)-x(T)
 \leq v(N)+v(C)-v(T).
\]
Moreover,
\[
 v(N)+v(C)-v(T)
 =\phi(S)+\bigl(\phi(T)-v(T)\bigr)
              -\bigl(\phi(C)-v(C)\bigr)
 \leq\phi(S).
\]
Exactness makes $F_C$ nonempty, so this face occurs in the minimum defining
$d_v^*(S)$ and proves the claim.
\end{proof}

The carrier comparison has a useful normal form that removes both slack
expressions.  It also converts failure of every candidate into a single
finite-dimensional separation problem.

\begin{proposition}[carrier marginal and face separation]
\label{prop:lx-carrier-separation}
Let $v$ be exact, $\phi=\phi(v)$, $S$ proper, and $P=N\setminus S$.
For every $k\in S$,
\[
 \begin{split}
 &[\phi(P\cup\{k\})-v(P\cup\{k\})]
   -[\phi_k-v(\{k\})]\\
 &\hspace{5em}=\phi(P)-[v(P\cup\{k\})-v(\{k\})].
 \end{split}
\]
Consequently $k$ satisfies the carrier inequality if and only if
\[
 v(P\cup\{k\})-v(\{k\})\geq\phi(P).
 \tag{M}\label{eq:lx-carrier-marginal}
\]
Let $K\subseteq S$ be nonempty.  If \eqref{eq:lx-carrier-marginal} fails
strictly for every $k\in K$, then every point of the exact face
$F_{P\cup\{k\}}$ satisfies $x(S)>\phi(S)$.  If, in addition,
$(\star_S)$ fails, there is a nonzero $w\in\R^S_+$ such that
\[
 \max_{k\in K}\ \max_{x\in F_{P\cup\{k\}}}
       \sum_{i\in S}w_i(x_i-\phi_i)<0.
 \tag{Sep}\label{eq:lx-carrier-separator}
\]
\end{proposition}

\begin{proof}
Because $P$ and $\{k\}$ are disjoint,
$\phi(P\cup\{k\})=\phi(P)+\phi_k$, which proves the identity and
\eqref{eq:lx-carrier-marginal}.  Now suppose the latter fails for $k$ and
take $x\in F_{P\cup\{k\}}$.  Core feasibility at $\{k\}$ gives
\[
 x(P)=v(P\cup\{k\})-x_k
 \leq v(P\cup\{k\})-v(\{k\})<\phi(P).
\]
Efficiency therefore gives $x(S)>\phi(S)$.

Let $Z$ be the convex hull of the union of the faces
$F_{P\cup\{k\}}$, $k\in K$.  It is a compact subset of the core, and every
point of $Z$ has $S$-sum strictly larger than $\phi(S)$.  Hence failure of
$(\star_S)$ implies that $Z$ is disjoint from
$\phi_S+\R^S_+$ after projection onto the $S$-coordinates.  Strong
separation gives a strict separating normal $w$.  Its coordinates must be
nonnegative, since otherwise its infimum on
$\phi_S+\R^S_+$ would be $-\infty$.  Translating by $\phi_S$ gives
\eqref{eq:lx-carrier-separator}.
\end{proof}

The separator is in fact forced by the star-failure cover; it is not an
independent normal that can be eliminated against that cover.

\begin{proposition}[residual-gap separator]
\label{prop:lx-residual-gap-separator}
Suppose failure of $(\star_S)$ has a tight-coalition certificate
\[
 r^S=\mu\mathbf1_N-\sum_T\lambda_T\mathbf1_T,
 \qquad \lambda_T\geq0,\quad
 \phi(T)=v(T)\ \text{if }\lambda_T>0,
\]
where $r_i\geq1$ on $S$.  Put
\[
 r_0=\min_{i\in S}r_i,\qquad
 K=\{i\in S:r_i=r_0\},\qquad
 H=S\setminus K,\qquad
 D=\sum_{i\in H}(r_i-r_0).
\]
Assume $H\neq\emptyset$.  If the carrier marginal fails at
$k\in K$ by
\[
 \gamma_k=[\phi(P\cup\{k\})-v(P\cup\{k\})]
           -[\phi_k-v(\{k\})]>0,
\]
then, for every $x\in F_{P\cup\{k\}}$,
\[
 \sum_{i\in H}\frac{r_i-r_0}{D}(x_i-\phi_i)
 \leq-\frac{r_0}{D}\gamma_k<0.
 \tag{RG}\label{eq:lx-residual-gap-separator}
\]
Consequently the normalized residual gap is a common strict separator for
every failed candidate in $K$, supported entirely on the higher-residual
class $H$.
\end{proposition}

\begin{proof}
For any $x\in\core(v)$ put $d=x-\phi$.  Efficiency gives $d(N)=0$, and
core feasibility plus tightness gives $d(T)\geq0$ whenever
$\lambda_T>0$.  Taking the scalar product of the cover identity with $d$
therefore yields
\[
 \sum_{i\in S}r_i d_i=-\sum_T\lambda_Td(T)\leq0.
\]
Splitting off the minimum residual gives
\[
 r_0d(S)+\sum_{i\in H}(r_i-r_0)d_i\leq0.
\]
For $x\in F_{P\cup\{k\}}$, the carrier calculation in
Proposition~\ref{prop:lx-carrier-separation} gives
$d(S)\geq\gamma_k>0$.  Substitution and division by $D$ prove
\eqref{eq:lx-residual-gap-separator}.
\end{proof}

\begin{corollary}[residual-gap escape criterion]
\label{cor:lx-residual-gap-escape}
Under the hypotheses of Proposition~\ref{prop:lx-residual-gap-separator}, if
some $k\in K$ satisfies
\[
 \max_{x\in F_{P\cup\{k\}}}
 \sum_{i\in H}\frac{r_i-r_0}{D}(x_i-\phi_i)\geq0,
 \tag{E}\label{eq:lx-residual-gap-escape}
\]
then $k$ satisfies the carrier inequality
\[
 [\phi(P\cup\{k\})-v(P\cup\{k\})]
 \leq[\phi_k-v(\{k\})].
\]
Thus existence of a minimum-residual carrier is reduced to escape from one
explicit scalar separator; no coordinatewise choice of a member of $K$ is
required.
\end{corollary}

\begin{proof}
The contrapositive is exactly inequality~\eqref{eq:lx-residual-gap-separator}.
\end{proof}

The strict separator also has a finite dual form.  This removes the
quantification over all points of the carrier-tight faces and exposes the
precise slack identity induced by the residual cover.

\begin{proposition}[dual form of the carrier-face separator]
\label{prop:lx-carrier-separator-dual}
Use the hypotheses and notation of
Proposition~\ref{prop:lx-carrier-separation}, and put
$T_k=P\cup\{k\}$.  A vector $w\in\R^S_+$ satisfies
\eqref{eq:lx-carrier-separator} if and only if, for every $k\in K$, there
are multipliers $\theta^k_U\geq0$ on the nonempty proper coalitions and
$a_k,b_k\in\R$ such that
\[
 w^S=a_k\mathbf1_N+b_k\mathbf1_{T_k}
          -\sum_{\emptyset\ne U\subsetneq N}
             \theta^k_U\mathbf1_U,
 \tag{D1}\label{eq:lx-separator-stationarity}
\]
where $w^S$ is $w$ on $S$ and zero on $P$, and
\[
 \sum_U\theta^k_U\,[\phi(U)-v(U)]
   < b_k[\phi(T_k)-v(T_k)].
 \tag{D2}\label{eq:lx-separator-slack}
\]
Whenever the carrier marginal fails at $k$, the right-hand slack in
\eqref{eq:lx-separator-slack} is positive and every such certificate has
$b_k>0$.
\end{proposition}

\begin{proof}
Fix $k$ and maximize $w^S\!\cdot x$ over the core subject to
$x(T_k)=v(T_k)$.  In the convention $-x(U)\leq-v(U)$ for the core
inequalities, the linear-programming dual has nonnegative multipliers
$\theta^k_U$, free multipliers $a_k,b_k$ on efficiency and the face
equality, and stationarity exactly
\eqref{eq:lx-separator-stationarity}.  Its objective is
\[
 -\sum_U\theta^k_Uv(U)+a_kv(N)+b_kv(T_k).
\]
Taking the scalar product of
\eqref{eq:lx-separator-stationarity} with $\phi$ shows that this objective
is strictly below $w\cdot\phi_S$ exactly when
\eqref{eq:lx-separator-slack} holds.  Strong duality therefore proves the
equivalence, simultaneously for all $k$ because the same $w$ occurs in
each stationarity identity.

If the carrier marginal fails, then
$\phi(T_k)-v(T_k)>\phi_k-v(\{k\})\geq0$.  The left side of
\eqref{eq:lx-separator-slack} is nonnegative, so strictness forces
$b_k>0$.
\end{proof}

\begin{corollary}[aggregated separator identity]
\label{cor:lx-aggregated-separator}
If the carrier marginal fails for every $k\in K$, then the separator in
Proposition~\ref{prop:lx-carrier-separation} supplies
$c>0$, $A\in\R$, a collection $\eta_U\geq0$, and a nonzero
$w\in\R^S_+$ such that
\[
 \mathbf1_K
 =c w^S-A\mathbf1_N-|K|\mathbf1_P
       +\sum_U\eta_U\mathbf1_U,
 \tag{AD1}\label{eq:lx-aggregate-separator}
\]
and
\[
 \sum_U\eta_U[\phi(U)-v(U)]
   <\sum_{k\in K}[\phi(P\cup\{k\})-v(P\cup\{k\})].
 \tag{AD2}\label{eq:lx-aggregate-separator-slack}
\]
More precisely, one may take
\[
 c=\sum_{k\in K}\frac1{b_k},\qquad
 A=\sum_{k\in K}\frac{a_k}{b_k},\qquad
 \eta_U=\sum_{k\in K}\frac{\theta_U^k}{b_k}.
\]
\end{corollary}

\begin{proof}
Proposition~\ref{prop:lx-carrier-separator-dual} gives $b_k>0$ for every
$k\in K$.  Divide \eqref{eq:lx-separator-stationarity} and
\eqref{eq:lx-separator-slack} by $b_k$ and sum over $k$.  Since
$\mathbf1_{P\cup\{k\}}=\mathbf1_P+e_k$ and
$\sum_{k\in K}e_k=\mathbf1_K$, the stationarity sum rearranges to
\eqref{eq:lx-aggregate-separator}; the strict slack inequalities sum to
\eqref{eq:lx-aggregate-separator-slack}.
\end{proof}

The same comparison has a quantitative converse at the level of total
$S$-payoff.  For a core point $\psi$, define
\begin{align*}
 \gamma_S(\psi)
  &=\min_{\emptyset\ne C\subsetneq S}
    \{[\psi(P\cup C)-v(P\cup C)]-[\psi(C)-v(C)]\},\\
 \eta_S(\psi)
  &=\max\{x(S)-\psi(S):x\in\core(v),\ x_i\geq\psi_i\ (i\in S)\}.
\end{align*}
Thus $(\star_S)$ at $\psi$ is equivalent to $\eta_S(\psi)>0$.

\begin{lemma}[exact carrier lift]
\label{lem:lx-carrier-lift}
If $v$ is exact, then for every nonempty proper $C\subset S$ there is a
$z^C\in\core(v)$ such that
\[
 z^C(S)-\psi(S)\geq
 [\psi(P\cup C)-v(P\cup C)]-[\psi(C)-v(C)].
\]
\end{lemma}

\begin{proof}
Choose $z^C$ tight at $P\cup C$.  Core feasibility at $C$ gives
$z^C(P)\leq v(P\cup C)-v(C)$.  Efficiency and the identity
$\psi(P)=v(N)-\psi(S)$ now give the displayed inequality.
\end{proof}

The only possible obstruction is coordinatewise: the carrier-tight points
raise the total payoff of $S$, but different members of $S$ may fall at
different points.  For a pair, convexity removes that obstruction without
loss.

There is a useful normalization that isolates exactly what remains to be
proved and, unexpectedly, removes the leximin selection from the statement.

\begin{proposition}[exact-core closure formulation]
\label{prop:lx-exact-core-closure}
Let $v$ be exact, let $\psi\in\core(v)$, and put
$D=\core(v)-\psi$ and $s_\psi(A)=\psi(A)-v(A)$.  Then
\[
 s_\psi(A)=-\min_{d\in D}d(A)=\max_{d\in D}d(N\setminus A).
 \tag{EC}\label{eq:lx-exact-core-support}
\]
Consequently
\[
 \gamma_S(\psi)=
 \min_{\emptyset\ne C\subsetneq S}
 \left\{\max_{d\in D}d(S\setminus C)+\min_{d\in D}d(C)\right\},
 \tag{CG}\label{eq:lx-convex-gamma}
\]
whereas
\[
 \eta_S(\psi)=\max\{d(S):d\in D,\ d_i\geq0\ (i\in S)\}.
 \tag{CE}\label{eq:lx-convex-eta}
\]
Thus the quantitative carrier inequality is a statement about every
translated exact core, not specifically about a leximin point.
\end{proposition}

\begin{proof}
Exactness gives $v(A)=\min_{x\in\core(v)}x(A)$, hence
$s_\psi(A)=-\min_{d\in D}d(A)$.  Every $d\in D$ has $d(N)=0$, so the last
quantity equals $\max_{d\in D}d(N\setminus A)$.  Substituting
$A=P\cup C$ and then $A=C$ gives \eqref{eq:lx-convex-gamma}; translating
the defining program for $\eta_S$ gives \eqref{eq:lx-convex-eta}.
\end{proof}

The closure qualification in this formulation is essential.  For the convex
hull in $\mathbb R^3$ of
\[
 (0,0,0),(4,2,-12),(-12,14,3),(18,9,-17),
 (10,17,-9),(4,-2,-5),(15,20,-5),
\]
the right side of \eqref{eq:lx-convex-gamma} is $3$, while the largest
nonnegative total in the hull is $0$.  Append to each point a fourth
coordinate equal to minus its total and take the lower-envelope exact game.
Its full core is the coalitional closure of those generators: the same
reference point still has $\gamma_S=3$, but now $\eta_S=27$, attained at
$(10,17,0,-27)$.  Thus arbitrary convexity is false by a gap of $3$, while
exact-core closure repairs this candidate by a gap of $24$.

The remaining obstruction has an exact dual that is formally the slack-priced
version of a failure cover.  Define the common-coordinate margin
\[
 \kappa_S(\psi)=
 \max\{z:d\in D,\ d(S)\geq\gamma_S(\psi),\ d_i\geq z\ (i\in S)\}.
\]

\begin{proposition}[quantitative orthant dual]
\label{prop:lx-quantitative-orthant-dual}
The value $\kappa_S(\psi)$ equals
\[
 \min\left\{
   \sum_T\lambda_Ts_\psi(T)-q\gamma_S(\psi):
   \begin{array}{l}
   \lambda_T,q,w_i\geq0,\quad \sum_{i\in S}w_i=1,\\
   u\mathbf1_N-\sum_T\lambda_T\mathbf1_T
      =q\mathbf1_S+w^S
   \end{array}
 \right\}.
 \tag{OD}\label{eq:lx-quantitative-orthant-dual}
\]
Here $u$ is free and $w^S$ is $w$ on $S$ and zero on $P$.
Consequently (Q) is equivalent to $\kappa_S(\psi)\geq0$, while the stronger
uniform lift asks for $\kappa_S(\psi)\geq\gamma_S(\psi)/|S|$.
\end{proposition}

\begin{proof}
Write the translated core inequalities as
$-d(T)\leq s_\psi(T)$, the high-total constraint as
$-d(S)\leq-\gamma_S(\psi)$, and the coordinate constraints as
$-d_i+z\leq0$.  The nonnegative dual multipliers are respectively
$\lambda_T,q,w_i$; efficiency has free multiplier $u$.  Stationarity in
$z$ gives $\sum_iw_i=1$, stationarity in $d$ gives the displayed incidence
identity, and the dual objective is
$\sum_T\lambda_Ts_\psi(T)-q\gamma_S(\psi)$.  Strong LP duality proves
\eqref{eq:lx-quantitative-orthant-dual}.
\end{proof}

\begin{proposition}[two-level orthant separators are impossible]
\label{prop:lx-two-level-orthant-separator}
Let $v$ be exact, $\psi\in\core(v)$, and $\gamma_S(\psi)>0$.  For every
nonempty $A\subseteq S$ there is a $d^A\in D$ such that
\[
 d^A(S)\geq\gamma_S(\psi),\qquad
 d^A(A)\geq\gamma_S(\psi).
 \tag{OI}\label{eq:lx-indicator-witness}
\]
Consequently, if $w=a\mathbf1_S+b\mathbf1_A$ with $a,b\geq0$, then
some $d\in D$ with $d(S)\geq\gamma_S(\psi)$ satisfies
$w\cdot d_S\geq(a+b)\gamma_S(\psi)$.
In particular, no strictly positive vector with at most two coordinate levels
can strictly separate
\[
 \{d_S:d\in D,\ d(S)\geq\gamma_S(\psi)\}
 \quad\hbox{from}\quad\mathbb R^S_+.
\]
Thus any failure of the quantitative carrier inequality admits a strictly
positive separator, but every such separator must have at least three distinct
coordinate levels.
\end{proposition}

\begin{proof}
If $A\subsetneq S$, put $C=S\setminus A$ and choose a core point tight at
$P\cup C$.  For its displacement $d^A$ from $\psi$,
\eqref{eq:lx-exact-core-support} gives
\[
 d^A(A)=s_\psi(P\cup C)
 \geq\gamma_S(\psi)+s_\psi(C)
 \geq\gamma_S(\psi),
\]
while Lemma~\ref{lem:lx-carrier-lift} gives
$d^A(S)\geq\gamma_S(\psi)$.  If $A=S$, the latter lemma with any nonempty
proper $C\subset S$ supplies the required point.  Hence
$w\cdot d^A_S=a d^A(S)+b d^A(A)
\geq(a+b)\gamma_S(\psi)>0$ whenever $w\ne0$.

If the displayed convex set and the nonnegative orthant are disjoint, strong
separation supplies a nonzero $w\geq0$ whose maximum on the first set is
strictly negative.  Since that set is compact and every one of its points has
$S$-sum at least the positive number $\gamma_S(\psi)$, adding a sufficiently
small multiple of $\mathbf1_S$ makes $w$ strictly positive while preserving
strict separation.  A strictly positive vector with at most two levels has
the form $a\mathbf1_S+b\mathbf1_A$, contradicting the first part.
\end{proof}

\begin{proposition}[codimension-one uniform carrier lift]
\label{prop:lx-codimension-one-uniform-lift}
Let $v$ be exact, let $\psi\in\core(v)$, and let
$P=N\setminus S=\{p\}$ be a singleton.  Put $m=|S|$.  Then
\[
 x_i=\psi_i+\frac{\gamma_S(\psi)}m\quad(i\in S),
 \qquad
 x_p=\psi_p-\gamma_S(\psi)
 \tag{UL}\label{eq:lx-uniform-carrier-lift}
\]
defines a point of $\core(v)$ whenever $\gamma_S(\psi)>0$.  Consequently
\[
 \eta_S(\psi)\geq\gamma_S(\psi),
\]
and every member of $S$ can be raised simultaneously by
$\gamma_S(\psi)/|S|$.
\end{proposition}

\begin{proof}
Write $t=\gamma_S(\psi)/m$ and $d=x-\psi$.  Efficiency is immediate.
If $T\subseteq S$, then $d(T)=t|T|\geq0\geq-s_\psi(T)$.

Now let $T=P\cup C$ with $C\subsetneq S$ and put $A=S\setminus C$.
If $C$ is nonempty, the definition of $\gamma_S$ gives
\[
 s_\psi(T)\geq\gamma_S(\psi)+s_\psi(C)
 \geq\gamma_S(\psi),
\]
and hence
\[
 d(T)=-d(A)=-t|A|\geq-\gamma_S(\psi)\geq-s_\psi(T).
\]
It remains only to check $T=P$.  Exactness makes $v$ superadditive on
disjoint coalitions: a core point tight at $P\cup C$ gives
$v(P\cup C)\geq v(P)+v(C)$.  Therefore
$s_\psi(P\cup C)\leq s_\psi(P)+s_\psi(C)$, and any nonempty proper
$C\subset S$ yields
\[
 \gamma_S(\psi)
 \leq s_\psi(P\cup C)-s_\psi(C)
 \leq s_\psi(P).
\]
Thus $d(P)=-\gamma_S(\psi)\geq-s_\psi(P)$.  Every coalition constraint
holds, so $x\in\core(v)$ and $x(S)-\psi(S)=mt=\gamma_S(\psi)$.
\end{proof}

\begin{proposition}[sharp quantitative pair bound and uniform split]
\label{prop:lx-quantitative-pair}
If $v$ is exact, $\psi\in\core(v)$, and $S$ is a proper two-player
coalition, then there is an $x\in\core(v)$ satisfying
\[
 x_i-\psi_i\geq\frac{\gamma_S(\psi)}2\quad(i\in S).
\]
In particular, $\eta_S(\psi)\geq\gamma_S(\psi)$, and the coefficient one
is sharp.
\end{proposition}

\begin{proof}
The claim is immediate if $\gamma_S(\psi)\leq0$, so write
$\gamma=\gamma_S(\psi)>0$.  Let
\[
 K=\{d_S:d\in D,\ d(S)\geq\gamma\}.
\]
If $K$ did not meet $(\gamma/2)\mathbf1_S+\mathbb R^S_+$, strong
separation would supply a nonzero $w\geq0$ such that
\[
 \max_{d_S\in K}w\cdot d_S
 <\frac\gamma2\sum_{i\in S}w_i.
 \tag{PS}\label{eq:lx-pair-separation}
\]
Every nonnegative two-dimensional $w$ can be written
$a\mathbf1_S+b\mathbf1_A$ for $a,b\geq0$ and a nonempty $A\subseteq S$.
Proposition~\ref{prop:lx-two-level-orthant-separator} gives a point of $K$
with weighted value at least $(a+b)\gamma$, whereas the right side of
\eqref{eq:lx-pair-separation} is
$\gamma(2a+b|A|)/2\leq(a+b)\gamma$, a contradiction.  Hence the uniform
split exists.  The exact lower-envelope equality fixtures show sharpness.
\end{proof}

\begin{corollary}[quantitative carrier bound through four players]
\label{cor:lx-quantitative-through-four}
For every exact game with $|N|\leq4$, every core point $\psi$, and every
proper nonsingleton $S$, one has
\[
 \eta_S(\psi)\geq\gamma_S(\psi).
\]
More generally, the inequality holds in every dimension whenever
$|S|=2$ or $|N\setminus S|=1$.  Hence any counterexample requires
$|S|\geq3$, $|N\setminus S|\geq2$, and therefore at least five players.
\end{corollary}

The first ostensibly open dimension can also be closed, although our present
proof is a finite exact computation.

\begin{proposition}[computer-assisted five-player carrier theorem]
\label{prop:lx-quantitative-through-five}
For every exact game with $|N|\leq5$, every core point $\psi$, and every
proper nonsingleton $S$,
\[
 \eta_S(\psi)\geq\gamma_S(\psi).
\]
In the sole new case $|N|=5$, $|S|=3$, the exactness assumption can be
weakened to total balancedness (nonemptiness of every subgame core).
\end{proposition}

\begin{proof}
The preceding results leave only $|N|=5$, $|S|=3$, and
$P=N\setminus S=\{p,q\}$.  The case $\gamma_S(\psi)\leq0$ is immediate.
Otherwise translate $\psi$ to zero, scale so that $\gamma_S(\psi)=1$, and
write $s(A)=s_\psi(A)$.  In fact, total balancedness suffices: a core point
of each subgame $T$, translated by $\psi|_T$ and extended arbitrarily outside
$T$ to have zero grand total, gives a vector $d^T$ such that
\[
 d^T(N)=0,\qquad d^T(T)=-s(T),\qquad
 d^T(U)\geq-s(U)\quad(\emptyset\ne U\subseteq T).
 \tag{X5}\label{eq:lx-five-exact-cone}
\]
Together with $s\geq0$, $s(\emptyset)=s(N)=0$, and the six carrier
inequalities
\[
 s(P\cup C)-s(C)\geq1
 \quad(\emptyset\ne C\subsetneq S),
 \tag{C5}\label{eq:lx-five-carrier-cone}
\]
these constraints define a rational polyhedron containing every normalized
totally balanced instance under consideration.

The LP defining $\eta_S$ has the slack-linear dual
\[
 \min_{\lambda,w,\rho\geq0}\sum_{\emptyset\ne T\subsetneq N}\lambda_Ts(T)
 \tag{E5}\label{eq:lx-five-eta-dual}
\]
subject to
\[
 \sum_{T\ni p}\lambda_T=\sum_{T\ni q}\lambda_T=1+\rho,
 \qquad
 \sum_{T\ni i}\lambda_T+w_i=\rho\quad(i\in S).
\]
This is a five-equation standard-form polyhedron with $34$ nonnegative
variables.  The verifier enumerates all $\binom{34}{5}=278{,}256$ candidate
bases by exact elimination.  They give $263$ distinct vertices, or $43$
orbits under the $S_3\times S_2$ player symmetry.  For each orbit
representative it then minimizes the objective in
\eqref{eq:lx-five-eta-dual} over
\eqref{eq:lx-five-exact-cone}--\eqref{eq:lx-five-carrier-cone}.
Independently replayed rational dual certificates give only the lower bounds
\[
 1,\quad \frac32,\quad \frac53,\quad 2,\quad 3,\quad 4;
\]
the minimum is $1$ and is attained.  Since a slack-nonnegative objective has
an optimal vertex in \eqref{eq:lx-five-eta-dual}, every normalized instance
has $\eta_S\geq1=\gamma_S$.  The exact enumeration and all $43$ rational
certificates are in
\path{src/audit_n5_exact_cone_eta_vertices.py} and
\path{results/n5_totally_balanced_cone_eta_vertices.json}.
\end{proof}

The exact-total uniform displacement used in the codimension-one proof does
not extend to this case.  An exact integer five-player game recorded in the
supplementary counterexample artifact has
$\gamma_S=1$, $\eta_S=2$, and $\kappa_S=1/2$, but no displacement with
$d(S)=1$ and $d_i=1/3$ for all $i\in S$.  Thus the five-player theorem
genuinely needs the freedom to gain more than $\gamma_S$ in total.

The eta dual has a finite combinatorial interpretation that removes the
apparently unbounded parameter $\rho$.

Minimal balanced collections and their generation are classical; complete
computational catalogs are now available through seven players
\citep{LaplaceMermoudEtAl2025}.  The point here is their target-conditioned
appearance as the vertices of the eta dual, followed by optimization over the
exact-game cone.  The latter cone has an independent facet description in
terms of indecomposable minimal semi-balanced systems
\citep{StudenyKratochvil2022}.

\begin{lemma}[eta-dual vertices are minimal balanced collections]
\label{lem:lx-eta-minimal-balanced}
Fix a proper nonsingleton $S\subset N$ and put $P=N\setminus S$.  On every
nonnegative slack system, the eta dual in \eqref{eq:lx-five-eta-dual} (with
the evident degree equation for every member of $P$) has an optimal vertex
with
\[
 \lambda_S=0,
 \qquad
 \lambda_{\{i\}}=0\quad(i\in S).
 \tag{DV}\label{eq:lx-dominated-eta-columns}
\]
After imposing \eqref{eq:lx-dominated-eta-columns}, its vertices are in
bijection with the minimal balanced collections on $N$ that contain $S$.
Explicitly, adjoining $S$ with coefficient one to a dual point gives
\[
 \mathbf1_S+
 \sum_T\lambda_T\mathbf1_T+
 \sum_{i\in S}w_i\mathbf1_{\{i\}}
   =(1+\rho)\mathbf1_N.
 \tag{MB}\label{eq:lx-eta-balanced-collection}
\]
Conversely, rescale a minimal balanced collection containing $S$ so that its
$S$-coefficient is one, interpret its $S$-singleton coefficients as the
$w_i$, and interpret all remaining coefficients except that of $S$ as the
$\lambda_T$.  In particular, every relevant dual vertex uses at most $|N|$
coalitions after $S$ is adjoined.
\end{lemma}

\begin{proof}
If $a=\lambda_S>0$, deleting $a$ from $\lambda_S$ and adding $a$ to every
$w_i$, $i\in S$, preserves all degree equations and weakly decreases the
objective by $a s(S)$.  Similarly, $\lambda_{\{i\}}$ can be moved into
$w_i$ at weakly lower cost.  Hence an optimum exists on the face
\eqref{eq:lx-dominated-eta-columns}, and a vertex optimum can be chosen there.

On that face the dual degree equations are exactly
\eqref{eq:lx-eta-balanced-collection}.  Thus the face is affinely isomorphic
to the slice, at $S$-coefficient one, of the nonnegative cone of balanced
collections containing $S$.  A point of this slice is a vertex precisely when
it lies on an extreme ray with positive $S$-coefficient: a decomposition into
two such rays gives a convex decomposition after normalization, while a
nonzero ray with zero $S$-coefficient gives a recession direction and again
precludes extremality.  The extreme rays of the balanced-collection cone are
the minimal balanced collections.  Their positive incidence columns are
linearly independent, so there are at most $|N|$ of them.  Finally, in the
reverse construction the normalized coefficient of $S$ is at most one
because every member of $S$ has total balanced degree one; hence
$\rho\geq0$ as required.
\end{proof}

The totally balanced strengthening is dimension-sharp.  Its first
six-player escape nevertheless exposes a reusable exactness identity.

\begin{lemma}[three-block exact-closure repair]
\label{lem:lx-three-block-exact-closure}
Let $N=\{1,\ldots,6\}$, $S=\{1,2,3\}$, $P=\{4,5,6\}$, and abbreviate
$s=s_\psi$.  For every exact game and every $\psi\in\core(v)$,
\[
 s(2345)+s(146)+s(56)\geq s(1456)-s(1).
 \tag{TB}\label{eq:lx-three-block-exact-closure}
\]
Consequently the eta-dual supported with unit weight on
$2345,146,56$ has cost at least $\gamma_S(\psi)$.
\end{lemma}

\begin{proof}
Translate $\psi$ to zero.  Choose core displacements $a,b,c$ tight at
$5,45,1456$, respectively.  The incidence identities
\begin{align*}
 \mathbf1_1+\mathbf1_{2345}+\mathbf1_{56}
   &=\mathbf1_N+\mathbf1_5,\\
 \mathbf1_5+\mathbf1_{2345}+\mathbf1_{146}
   &=\mathbf1_N+\mathbf1_{45},\\
 \mathbf1_1+\mathbf1_{45}+\mathbf1_{146}+\mathbf1_{56}
   &=2\,\mathbf1_{1456}
\end{align*}
and core feasibility give
\begin{align*}
 s(5)&\leq s(1)+s(2345)+s(56),\\
 s(45)&\leq s(5)+s(2345)+s(146),\\
 2s(1456)&\leq s(1)+s(45)+s(146)+s(56).
\end{align*}
Substituting the first two bounds into the third and dividing by two proves
\eqref{eq:lx-three-block-exact-closure}.  The final assertion follows from
$s(1456)-s(1)\geq\gamma_S(\psi)$, the carrier inequality for $C=\{1\}$.
\end{proof}

There is an integer totally balanced game for this split with
$\gamma_S=1$, $\eta_S=0$, and zero-cost eta dual $\{2345,146,56\}$.
It is not exact; its exact-core closure is the zero game, so its carrier gap
also collapses to zero.  Thus Lemma~\ref{lem:lx-three-block-exact-closure}
identifies exactly the crossing constraint supplied by exactness.  The game,
closure, and incidence proof are independently checked in the supplementary
three-block verifier.

The repair is not isolated: it closes the complete first multilevel layer of
the six-player eta dual.

\begin{proposition}[complete six-player unit-triple layer]
\label{prop:lx-n6-unit-triple-layer}
In the split $|S|=|P|=3$, consider eta-dual solutions with $\rho=1$,
$w=0$, and three unit-weight supporting coalitions.  Equivalently, every
target player has support degree one and every outside player has degree two.
For every such dual and every exact normalized slack system satisfying the
carrier inequalities, its slack cost is at least one.  Hence this entire
dual layer satisfies the quantitative carrier bound.
\end{proposition}

\begin{proof}
Each target player chooses its unique supporting coalition, while each
outside player chooses the unique coalition from which it is absent.
Enumerating these two maps and quotienting by $S_3\times S_3$ gives exactly
$16$ support orbits.  Minimization over the global exact-witness cone gives
exact values one for seven orbits and two for nine.  Every value is certified
by independently replayed rational primal and dual solutions.  The complete
enumeration is in \path{src/audit_n6_unit_triple_eta_duals.py}, with the
certificates in \path{results/n6_unit_triple_eta_duals.json}.
\end{proof}

The preceding layer is the sharp boundary of a complete six-player theorem.
The full census below was the discovery proof, but the separator geometry
reduces it to a much smaller rank classification.

\begin{lemma}[rank-six three-level classification]
\label{lem:lx-n6-three-level-classification}
Consider an undominated eta-dual vertex for one of the two unresolved
six-player splits.  If its target-floor multipliers $w_i$ have at least three
distinct values, then, up to permutations preserving $S$ and $P$, necessarily
$|S|=4$ and the associated normalized minimal balanced collection is
\[
 \mathbf1_S+e_1+2e_2+\mathbf1_{35}+\mathbf1_{136}
     +2\,\mathbf1_{456}=3\,\mathbf1_N.             \tag{TL}
 \label{eq:lx-n6-exceptional-three-level}
\]
In particular its eta objective is
\[
 s(35)+s(136)+2s(456)\geq2\gamma_S.
 \tag{TC}\label{eq:lx-n6-exceptional-cost}
\]
There is no three-level vertex in the $3+3$ split.
\end{lemma}

\begin{proof}
By Lemma~\ref{lem:lx-eta-minimal-balanced}, adjoining $S$ and the positive
target-singleton columns gives a minimal balanced collection with at most six
members.  Let $q$ be the number of positive $w_i$.  The collection cannot
contain $S$ and all target singletons, since their incidence columns are
linearly dependent.  Thus at least one $w_i$ is zero.  Three distinct levels
therefore require at least two distinct positive singleton coefficients, so
\[
 2\leq q\leq |S|-1.
\]
After $S$ and these $q$ singleton columns are fixed, the rank-six bound leaves
at most $5-q\leq3$ further coalition columns.  This is a small finite
$0$--$1$ incidence problem, not an exact-game cone optimization.  Exact
determinant elimination gives
\[
\begin{array}{c|r|r|r}
(|S|,|P|)&\text{supports checked}&\text{three-level supports}
 &\text{symmetry orbits}\\ \hline
(3,3)&32{,}567&0&0\\
(4,2)&32{,}566&8&1.
\end{array}
\]
The unique orbit is \eqref{eq:lx-n6-exceptional-three-level}.  These counts,
positivity of the balancing weights, column independence, and orbit uniqueness
are replayed in exact integer arithmetic by
\path{src/verify_n6_three_level_eta_duals.py}, with output
\path{results/n6_three_level_eta_duals.json}.

For the exceptional orbit, the carrier inequality with $C=\{4\}$ is
$s(456)-s(4)\geq\gamma_S$.  Slack nonnegativity therefore gives
\[
 s(35)+s(136)+2s(456)
 \geq2s(456)
 \geq2\gamma_S+2s(4)
 \geq2\gamma_S,
\]
which proves \eqref{eq:lx-n6-exceptional-cost}.
\end{proof}

\begin{proposition}[rank-reduced six-player carrier theorem]
\label{prop:lx-quantitative-through-six}
For every exact game with $|N|\leq6$, every core point $\psi$, and every
proper nonsingleton $S$,
\[
 \eta_S(\psi)\geq\gamma_S(\psi).
\]
\end{proposition}

\begin{proof}
Corollary~\ref{cor:lx-quantitative-through-four},
Proposition~\ref{prop:lx-quantitative-through-five}, and the general cases
$|S|=2$ and $|P|=1$ leave only the six-player splits
$(|S|,|P|)=(3,3)$ and $(4,2)$.  Suppose, to the contrary, that
$\gamma=\gamma_S(\psi)>0$ and $\eta_S(\psi)<\gamma$.  Translate $\psi$ to
zero and choose an optimal eta-dual vertex on the undominated face from
Lemma~\ref{lem:lx-eta-minimal-balanced}; denote its target-floor multipliers
by $w_i$ and its value by $L=\sum_T\lambda_Ts(T)=\eta_S(\psi)$.

Put
\[
 K=\{d_S:d\in D,\ d(S)\geq\gamma\}.
\]
The dual degree equations, dotted with any $d\in D$, give
\[
 d(S)+\sum_{i\in S}w_i d_i
 =-\sum_T\lambda_Td(T)
 \leq\sum_T\lambda_Ts(T)=L.
\]
Hence for every $d_S\in K$,
\[
 \sum_{i\in S}w_i d_i\leq L-d(S)<0.
\]
Thus $w\geq0$ strictly separates $K$ from the nonnegative orthant.  If $w$
had at most two coordinate levels, it would have the form
$a\mathbf1_S+b\mathbf1_A$ with $a,b\geq0$, contradicting
Proposition~\ref{prop:lx-two-level-orthant-separator}.  Therefore $w$ has at
least three levels.

Lemma~\ref{lem:lx-n6-three-level-classification} now excludes the $3+3$
split and leaves only the exceptional $4+2$ orbit, whose dual value is at
least $2\gamma$, again contradicting $L<\gamma$.  This proves the theorem.

For independent redundancy, the original exhaustive audit remains available.
Its C++ census checks $6{,}508{,}883$ supports per split, finds $21{,}695$ and
$19{,}501$ target-conditioned minimal balanced collections, and reduces them
to $870$ and $619$ symmetry orbits.  All $1{,}489$ exact-cone optima are at
least one by replayable rational certificates.  The C++ census is
\path{src/enumerate_n6_eta_dual_vertices.cpp}; its exact output and independent
structural verification are
\path{results/n6_eta_dual_vertex_census.json} and
\path{results/n6_eta_dual_vertex_census_verification.json}.  The latter
rechecks the support count, balancing equations, incidence rank, orbit
uniqueness, and recovery of all $505$ previously sampled orbits.  The exact
LP verifier and its full rational certificates are
\path{src/audit_n6_eta_dual_vertices.py} and
\path{results/n6_eta_dual_vertex_audit.json}.  Finally,
\path{src/verify_n6_eta_dual_vertex_audit.py} replays all $1{,}489$ saved
lower-bound identities in exact arithmetic without calling an optimizer; its
summary is \path{results/n6_eta_dual_vertex_audit_verification.json}.  None of
this larger audit is needed for the rank-reduced proof above.
\end{proof}

The six-player threshold is exact.

\begin{proposition}[sharp seven-player exact counterexample]
\label{prop:lx-quantitative-n7-counterexample}
There is an exact seven-player game whose leximin point $\phi$ and a proper
coalition $S$ satisfy
\[
 \gamma_S(\phi)=\frac34,
 \qquad
 \eta_S(\phi)=0.
\]
For the same game the core reference $\psi=0$ satisfies
$(\gamma_S(\psi),\eta_S(\psi))=(1,0)$.
Consequently the quantitative carrier inequality is false in unrestricted
dimension, and Proposition~\ref{prop:lx-quantitative-through-six} is sharp in
the number of players.
\end{proposition}

\begin{proof}
Let $N=\{1,\ldots,7\}$, $S=\{1,2,3\}$, $P=\{4,5,6,7\}$, and put
\[
 a=(-3,1,3,-1,2,-1,-1),
 \qquad
 b=(3,-2,0,-1,-1,-1,2).
\]
Both vectors have total zero.  Define the lower-envelope game
\[
 v(T)=\min\{0,a(T),b(T)\}
 \qquad(T\subseteq N),
\]
and take $\psi=0$.  Each of $0,a,b$ belongs to $\core(v)$, and at least one
of them is tight at every coalition.  Hence $v$ is exact and
$\psi\in\core(v)$.

Writing $s(T)=-v(T)$, direct evaluation gives
\[
\begin{array}{c|rrrrrr}
 C&1&2&3&12&13&23\\ \hline
 s(C)&3&2&0&2&0&2\\
 s(P\cup C)&4&3&1&3&1&3
\end{array}
\]
and therefore every one of the six carrier differences equals one.  Thus
$\gamma_S(0)=1$.

It remains to calculate $\eta_S(0)$.  The four coalitions
\[
 Q_1=13456,
 \quad Q_2=457,
 \quad Q_3=3467,
 \quad Q_4=567
\]
all have value zero, and their incidence vectors satisfy
\[
 \mathbf1_S+\sum_{j=1}^4\mathbf1_{Q_j}+e_1+2e_2
 =3\,\mathbf1_N.
 \tag{N7}\label{eq:lx-n7-sharp-identity}
\]
If $x\in\core(v)$ and $x_i\geq0$ for $i\in S$, then
$x(Q_j)\geq v(Q_j)=0$ for every $j$.  Taking the scalar product of
\eqref{eq:lx-n7-sharp-identity} with $x$ yields
\[
 x(S)+\sum_{j=1}^4x(Q_j)+x_1+2x_2=3x(N)=0.
\]
Every term on the left is nonnegative, so $x(S)=0$.  Since $x=0$ is
feasible, $\eta_S(0)=0$.

The leximin point can be certified in three elementary stages, without an
optimization oracle.  The singleton floors are
\[
 m=(-3,-2,0,-1,-1,-1,-1).
\]
At the first stage, $v(125)=0$ and efficiency imply $x(3467)\leq0$ for
every $x\in\core(v)$.  If all seven gains $x_i-m_i$ are at least $t$, then
\[
 x(3467)\geq m_3+m_4+m_6+m_7+4t=-3+4t,
\]
so $t\leq3/4$.  Equality fixes players $3,4,6,7$ at
$(3/4,-1/4,-1/4,-1/4)$.

On this first-stage face their coordinates sum to zero.  Since
$v(13)=v(27)=0$, the core constraints give $x_1\geq-3/4$ and
$x_2\geq1/4$.  Efficiency therefore gives $x_5\leq1/2$, so the second
common gain among the free players is at most $x_5-m_5\leq3/2$; equality
fixes player~5 at $1/2$.  Finally $x_1+x_2=-1/2$.  If their remaining two
gains are at least $u$, then
\[
 -\frac12=x_1+x_2\geq m_1+m_2+2u=-5+2u,
\]
so $u\leq9/4$.  All three bounds are attained by
\[
 \phi=\frac14a
 =\left(-\frac34,\frac14,\frac34,-\frac14,
          \frac12,-\frac14,-\frac14\right).
\]
Thus the successive Maschler strata are
\[
 \{3,4,6,7\}\text{ at }\frac34,\qquad
 \{5\}\text{ at }\frac32,\qquad
 \{1,2\}\text{ at }\frac94,
\]
which proves directly that this point is $\phi(v)$.
Direct evaluation of $s_\phi(T)=\phi(T)-v(T)$ gives
$s_\phi(P\cup C)-s_\phi(C)=3/4$ for all six nonempty proper
$C\subset S$, so $\gamma_S(\phi)=3/4$.  Moreover each $Q_j$ above is
$\phi$-tight.  If $x\in\core(v)$ dominates $\phi$ on $S$, put
$d=x-\phi$.  Taking the scalar product of
\eqref{eq:lx-n7-sharp-identity} with $d$ gives
\[
 d(S)+\sum_{j=1}^4d(Q_j)+d_1+2d_2=3d(N)=0.
\]
Every term on the left is nonnegative: the coordinate terms by domination,
and the coalition terms because $x(Q_j)\geq v(Q_j)=\phi(Q_j)$.  Hence
$d(S)=0$, while $x=\phi$ is feasible, proving $\eta_S(\phi)=0$.

The construction and all $128$ lower-envelope, core, carrier, degree, and
zero-cost identities are replayed in integer arithmetic by
\path{src/verify_n7_exact_carrier_counterexample.py}.  The exact rational
leximin calculation and the six $3/4$ differences are independently replayed
by \path{src/verify_n7_exact_carrier_counterexample_leximin.py}.
\end{proof}

The leximin example has four outside players.  At an arbitrary core reference,
three already suffice.

\begin{proposition}[exact outside-triple obstruction]
\label{prop:lx-n7-outside-triple}
There is an exact seven-player game, a core point $\psi$, and a target
$S$ with $|S|=4$ and $|N\setminus S|=3$ such that
\[
 \gamma_S(\psi)=1,
 \qquad
 \eta_S(\psi)=0.
\]
This example is not a leximin carrier counterexample: at its leximin point
all fourteen carrier differences are zero.
\end{proposition}

\begin{proof}
Let $S=1234$, $P=567$, and define
\[
 A=(-4,1,2,2,-3,1,1),
 \qquad
 B=(2,-1,0,0,1,-1,-1),
\]
and $v(T)=\min\{0,A(T),B(T)\}$.  As before, the three generators are
efficient core allocations and one is tight at every coalition, so $v$ is
exact and $\psi=0$ lies in its core.  Direct evaluation gives
\[
 s(P\cup C)-s(C)=1
 \qquad(\emptyset\ne C\subsetneq S),
\]
and hence $\gamma_S(0)=1$.

The four coalitions $235,456,457,1367$ have value zero, and
\[
 \mathbf1_S+e_1+3e_2+\mathbf1_{235}
 +2\,\mathbf1_{456}+2\,\mathbf1_{457}+3\,\mathbf1_{1367}
 =5\,\mathbf1_N.                                    \tag{OT}
 \label{eq:lx-n7-outside-triple-identity}
\]
If $x\in\core(v)$ and $x_i\geq0$ for $i\in S$, every term obtained by
pairing the left side with $x$ is nonnegative.  Their sum is $5x(N)=0$,
so $x(S)=0$.  The zero allocation is feasible, proving $\eta_S(0)=0$.
All lower-envelope values, the fourteen unit carrier differences, the
degree-five identity, and its rank-seven support are checked without an
optimizer by
\path{src/verify_n7_exact_carrier_outside_triple_counterexample.py}.
An accompanying exact-rational verifier gives the leximin point
\[
 (-2,1/3,4/3,4/3,-5/3,1/3,1/3)
\]
and zero for all fourteen carrier differences.
\end{proof}

The preceding construction is minimal in outside-set size among
seven-player obstructions at arbitrary core references.

\begin{proposition}[seven-player outside-pair theorem]
\label{prop:lx-n7-outside-pair}
For every exact seven-player game, every core point $\psi$, and every proper
nonsingleton $S$ with $|N\setminus S|\leq2$,
\[
 \eta_S(\psi)\geq\gamma_S(\psi).
\]
Consequently three is the minimum number of outside players in a
seven-player arbitrary-reference obstruction.  For leximin references, the
known minimum remains between three and four.
\end{proposition}

\begin{proof}
The codimension-one case is
Proposition~\ref{prop:lx-codimension-one-uniform-lift}, so only the split
$(|S|,|P|)=(5,2)$ remains.  If a violation existed, the proof of
Proposition~\ref{prop:lx-quantitative-through-six} would produce an
undominated eta-dual vertex whose target-floor multiplier has at least three
levels.  After adjoining $S$, its support is a minimal balanced collection
of at most seven columns.  If $q$ positive target-singleton columns occur,
three levels force $2\leq q\leq4$, leaving at most $6-q$ other coalition
columns.

Exact determinant enumeration of precisely these supports checks
$8{,}798{,}030$ candidates.  It finds $2{,}196$ three-level supports.  The
carrier inequalities and slack nonnegativity directly bound all but $96$ of
them; quotienting the latter by $S_5\times S_2$ leaves five orbits.  For each
orbit, minimization over the global exact-witness slack cone has exact value
$3$, with independently replayable rational dual certificates.  Hence no
three-level vertex has cost below the normalization
$\gamma_S(\psi)=1$, a contradiction.

For comparison, the same focused census in the $4+3$ split checks
$9{,}086{,}132$ supports and leaves $34$ orbits after the direct carrier
bound.  Their exact-value distribution is
\[
 14\text{ at }0,
 \qquad 12\text{ at }1,
 \qquad 8\text{ at }2.
\]
Thus the outside-triple example reflects several inequivalent rank-seven
obstruction mechanisms rather than a numerical accident.  The
accompanying integer census, exact cone audit, and optimizer-free replay are
all listed in the supplementary artifact; the replay checks every saved
rational lower-bound identity without calling an optimizer.
Proposition~\ref{prop:lx-n7-outside-triple} supplies the matching failure at
outside size three.
\end{proof}

The sharp counterexample propagates to every higher dimension.

\begin{corollary}[complete dimension threshold]
\label{cor:lx-quantitative-complete-threshold}
The quantitative carrier inequality holds for every exact game if
$|N|\leq6$, and fails for an exact game at every dimension $|N|\geq7$.
\end{corollary}

\begin{proof}
The positive statement is
Proposition~\ref{prop:lx-quantitative-through-six}.  Starting from the game
in Proposition~\ref{prop:lx-quantitative-n7-counterexample}, append any
number of null players: extend each of the three generators by zero and put
$v^+(T)=v(T\cap N)$.  This remains an exact lower-envelope game.

Its core is exactly $\core(v)\times\{0\}^k$.  Indeed, singleton rationality
makes each new coordinate nonnegative, while the old-grand-coalition
constraint and efficiency make their sum nonpositive; hence all are zero.
The converse inclusion is immediate from $v^+(T)=v(T\cap N)$.  The new
players therefore contribute fixed zero gains before the three positive
leximin strata, so the old leximin coordinates are unchanged.  For the same
target $S$, the enlarged outside set merely adds null players, and hence all
six carrier differences remain $3/4$.  The feasible core displacements on
$S$ are also unchanged, so $\eta_S=0$.
The identities are replayed through arbitrary prescribed finite dimension by
\path{src/verify_n7_exact_carrier_dummy_extensions.py}.
\end{proof}

The example has a rigid two-dimensional core but is not an isolated game.

\begin{proposition}[triangular core and generator minimality]
\label{prop:lx-n7-triangular-core}
For the seven-player game in
Proposition~\ref{prop:lx-quantitative-n7-counterexample},
\[
 \core(v)=\operatorname{conv}\{0,a,b\}.
\]
Moreover every representation of $v$ as the lower envelope of core
allocations requires at least three allocations.  Thus the displayed
three-generator representation is minimal.
\end{proposition}

\begin{proof}
The eight coalitions
\[
 \mathcal R=\{1234,125,1236,13456,27,457,3467,567\}
\]
all have value zero and satisfy
$\sum_{R\in\mathcal R}\mathbf1_R=4\,\mathbf1_N$.  Hence every core point is
tight on every member of $\mathcal R$: all eight core inequalities are
nonnegative, while their sum is four times efficiency.  Their incidence
matrix has rank five.  Both $a$ and $b$ lie in its kernel and are linearly
independent, so every $x\in\core(v)$ has the form $x=\alpha a+\beta b$.

The core inequalities at $3$, $13$, and $4$ now read, respectively,
$3\alpha\geq0$, $3\beta\geq0$, and
$-\alpha-\beta\geq-1$.  Thus
$\alpha,\beta\geq0$ and $\alpha+\beta\leq1$.  The reverse inclusion holds
because $0,a,b\in\core(v)$ and the core is convex.

Finally, among the three vertices, $a$ uniquely attains $v(1)$, $b$ uniquely
attains $v(2)$, and zero uniquely attains $v(123)$.  Any lower-envelope
representation by core allocations must therefore contain all three.
The rank, degree, core, and unique-attainment calculations are replayed by
\path{src/verify_n7_exact_carrier_structure.py}.
\end{proof}

\begin{proposition}[nonisolated two-parameter family]
\label{prop:lx-n7-carrier-family}
Let
\[
 \kappa\geq0,
 \qquad \theta-2\kappa>\frac14,
 \qquad \kappa+\theta\leq1.
\]
Then the exact game
\[
 v_{\kappa,\theta}(T)
 =\min\{0,a(T),(\kappa a+\theta b)(T)\}
\]
has leximin point $\phi=a/4$ and, for $S=123$,
\[
 \gamma_S(\phi)=\theta-2\kappa-\frac14>0,
 \qquad \eta_S(\phi)=0.
\]
Consequently the seven-player counterexample is neither isolated nor unique
up to permutation and common positive scaling.
\end{proposition}

\begin{proof}
Exactness again follows from the equal-total lower-envelope construction, and
$a/4$ is a core point.  The singleton floors are
\[
 (-3,\kappa-2\theta,0,-1,2\kappa-\theta,-1,-1).
\]
The first leximin stage is unchanged: $v_\theta(125)=0$ forces
$x(3467)\leq0$, so the first common gain is at most $3/4$ and fixes
$3,4,6,7$ at the coordinates of $a/4$.  On that face the constraints at
$13$ and $27$ give $x_1\geq-3/4$ and $x_2\geq1/4$, whence efficiency gives
$x_5\leq1/2$.  The next gain is therefore at most
$1/2-2\kappa+\theta$, attained by $a/4$.  It is strictly below the gain of
player~2 because $\kappa+\theta>1/4$, and below the gain of player~1 because
$\kappa+\theta\leq1$.

After fixing player~5, efficiency gives $x_1+x_2=-1/2$.  The preceding
bound on $x_1$ gives $x_2\leq1/4$, so the next gain is at most
$1/4-\kappa+2\theta$, again attained by $a/4$ and at most $9/4$; the
remaining coordinate is then fixed.  This proves the leximin assertion
directly.  Evaluation of the six carrier differences gives $3/4$ for the
three $C$ containing player~1, $\kappa+\theta-1/4$ for $C=2,23$, and
$\theta-2\kappa-1/4$ for $C=3$.  The last quantity is the minimum under the
parameter restrictions.  The four coalitions in
\eqref{eq:lx-n7-sharp-identity} remain zero-valued and $\phi$-tight, so the
same identity proves $\eta_S(\phi)=0$.

The region contains a relatively open two-dimensional wedge (take
$\kappa>0$ and both remaining inequalities strict), giving genuine
two-parameter robustness.  The slice $\kappa=0$ includes the earlier
$\theta=1$ and $\theta=4/5$ examples, whose singleton-floor multisets are not
proportional under any permutation.  Exact family checks are in
\path{src/verify_n7_exact_carrier_family.py}.
\end{proof}

\begin{corollary}[the carrier failure is not bumpable]
\label{cor:lx-n7-carrier-unbumpable}
For every game in Proposition~\ref{prop:lx-n7-carrier-family}, the exact bump
capacity at $S=123$ is zero:
\[
 d^*_{v_\theta}(123)=0.
\]
Hence this sharp failure of the quantitative carrier inequality cannot itself
produce a counterexample to proper-coalition monotonicity.
\end{corollary}

\begin{proof}
This is forced by the same zero-cost face geometry, not merely by unique
attainment at one vertex.  Three of the four coalitions in
\eqref{eq:lx-n7-sharp-identity} satisfy the incidence identity
\[
 \mathbf1_{135}-2\,\mathbf1_{123}
 =\mathbf1_{13456}+\mathbf1_{457}+\mathbf1_{567}
   -2\,\mathbf1_N.                                    \tag{NB}
\]
All three coalitions on the right have value zero throughout the family.
Hence every core point $x$ obeys
\[
 x(135)-2x(123)
 =x(13456)+x(457)+x(567)\geq0.
\]
The coalition $135$ also has value zero (its three generator payoffs are
$0,2,2(\kappa+\theta)$).  Thus every $x\in F_{135}$ has $x(123)\leq0$, while core
feasibility at $123$ gives the reverse inequality.  Consequently
\[
 \max_{x\in F_{135}}x(123)-v_\theta(123)=0.
\]
This term occurs in the minimum defining $d^*_{v_\theta}(123)$, while bump
capacity is nonnegative, proving equality.  Thus no positive exact bump at
this target exists.  In particular, the unbumpability is intrinsic to the
tight balanced certificate: any deformation preserving these three
zero-valued faces and $v(123)=v(135)=0$ remains blocked.
\end{proof}

The negative mechanism also exposes the rank threshold.  After adjoining
$S$, the zero-cost eta support in \eqref{eq:lx-n7-sharp-identity} is a minimal
balanced collection of seven linearly independent incidence columns.  The
outside-triple identity \eqref{eq:lx-n7-outside-triple-identity} has the same
rank-seven feature.  Six players permit at most six independent columns;
combined with the two-level separator theorem, this leaves only the single
rank-six three-level orbit in
\eqref{eq:lx-n6-exceptional-three-level}, which one carrier inequality
defeats.  Seven players are the first dimension in which multiple residual
three-level orbits survive that rank squeeze, and fourteen of them have zero
cost on the exact-game cone.  This is the conceptual six/seven transition;
the larger $1{,}489$-orbit audit is only independent redundancy.

The slack certificate is automatic on the most important thin-core model
class used by the searches.

\begin{proposition}[two-generator lower envelopes]
\label{prop:lx-two-generator}
Let $p,q\in\mathbb R^N$ have the same total, and define the exact game
\[
 v(T)=\min\{p(T),q(T)\}\quad(T\subsetneq N),
 \qquad v(N)=p(N)=q(N).
\]
Let $S$ be proper.  If $\max_{x\in\core(v)}x(S)>\phi(S)$ and
$(\star_S)$ fails at $\phi$, then there is a nonempty proper $C\subset S$
satisfying the hypothesis of
Proposition~\ref{prop:lx-carrier-capacity}.  Consequently
$d_v^*(S)\leq\phi(S)-v(S)$.
\end{proposition}

\begin{proof}
Both $p$ and $q$ belong to $\core(v)$, and one of them attains $v(T)$ at
every proper $T$, so $v$ is exact.  Put $P=N\setminus S$.  Exactness at $P$ gives
\[
 \max_{x\in\core(v)}x(S)=v(N)-v(P)=\max\{p(S),q(S)\}.
\]
After interchanging $p$ and $q$, binding therefore gives
$p(S)>\phi(S)$.  The two generator sums on $S$ cannot tie, since then
$v(S)=\max_{x\in\core(v)}x(S)$; hence $p(S)>q(S)$ and $p(P)<q(P)$.  Define
\[
 C=\{i\in S:p_i>\phi_i\}.
\]
This set is nonempty because $p(S)>\phi(S)$, and it is proper because
otherwise $p\in\core(v)$ would itself witness $(\star_S)$.  For $i\in C$,
core feasibility at the singleton gives
\[
 p_i>\phi_i\geq v(\{i\})=\min\{p_i,q_i\},
\]
so $q_i\leq\phi_i<p_i$.  Hence
\[
 p(C)-q(C)
 \geq p(C)-\phi(C)
 \geq p(S)-\phi(S)
 =\phi(P)-p(P).
\]
In particular $q(C)<p(C)$, and therefore
\[
 \begin{aligned}
 v(P\cup C)-v(C)
 &=\min\{p(P)+p(C),q(P)+q(C)\}-q(C)\\
 &=\min\{p(P)+p(C)-q(C),q(P)\}\\
 &\geq\phi(P).
 \end{aligned}
\]
The last inequality also uses
$\phi(P)\leq v(N)-v(S)=q(P)$, since $q(S)<p(S)$.
Equivalently,
$\phi(P\cup C)-v(P\cup C)\leq\phi(C)-v(C)$, which is precisely the
off-side slack-domination hypothesis.
\end{proof}

\begin{lemma}[exact-face compensation identity]
\label{lem:lx-generator-compensation}
Let $v$ be exact, let $P=N\setminus S$, and let $\psi\in\core(v)$.  For
nonempty $C\subsetneq S$, choose any $a\in\core(v)$ with
$a(P\cup C)=v(P\cup C)$; such an $a$ exists by exactness.  Then
\[
 v(P\cup C)-v(C)-\psi(P)
 =[a(P)-\psi(P)]+[a(C)-v(C)].
 \tag{GC}\label{eq:lx-generator-compensation}
\]
In particular, $a(P)\geq\psi(P)$ is sufficient for the off-side carrier
inequality at $C$; otherwise the only missing quantity is the nonnegative
trace compensation $a(C)-v(C)$.  Thus the identity is not restricted to a
finite lower-envelope representation.
\end{lemma}

\begin{proof}
Since $a(P\cup C)=v(P\cup C)$ and $P\cap C=\varnothing$, the identity is
algebraic.  Core feasibility gives $a(C)-v(C)\geq0$.
\end{proof}

So the bump is realized by a continuous path of \emph{exact} games in which only
the lower bound $\tau$ on $x(S)$ rises. This isolates the role of exactness: if
endpoint exactness is weakened to balancedness, the analogous interior path
need not be exact, and that is precisely the room inhabited by the
counterexamples of Section~\ref{sec:lx-necessity}. For proper nonsingleton
bumps on exact games, coalitional monotonicity is therefore equivalent to the
one-parameter statement:

\begin{lemma}[monotone path --- \textsc{open}]\label{lem:lx-C}
Along the exact path $v_\tau$, the map $\tau\mapsto\phi_i(v_\tau)$ is nondecreasing
for every $i\in S$ (and, empirically, nonincreasing for $i\notin S$).
\end{lemma}

Two reductions cut the work down further. First, the common case is trivial:

\begin{lemma}[non-binding]\label{lem:lx-nonbinding}
If $\phi(v)(S)\ge v(S)+d$ then $\phi(v')=\phi(v)$.
\end{lemma}
\begin{proof}
$\phi(v)$ is the leximin point of $\core(v)\supseteq\core(v')$ and lies in
$\core(v')$, hence is the leximin point of the smaller set too; $\phi$ is
single-valued.
\end{proof}

\noindent Binding ($\phi(v)(S)<v(S)+d$) requires $\phi(v)(S)<\max_{\core}x(S)$ and
is empirically rare. In every binding case observed (thousands of exact-preserving
bumps at $n\le6$, zero violations) one finds $\phi(v')(S)=v(S)+d$ and the two-sided
pattern $\phi_i(v')\ge\phi_i(v)$ for $i\in S$, $\le$ for $i\notin S$.

\paragraph{A cone criterion.} For $x=\phi(w)$ with tight set
$\mathcal T(x)=\{T:x(T)=w(T)\}$, say $(\star_S)$ holds if there is $e\in\R^N$ with
\[
  e(N)=0,\quad e(S)>0,\quad e_i\ge0\ \ (i\in S),\quad e(T)\ge0\ \ (T\in\mathcal T(x)):
\]
a first-order direction raising $x(S)$ without lowering any member of $S$ or
breaking a tight core constraint. Empirically $(\star_S)$ is \emph{necessary} for
the bump to be monotone (it fails at $\phi(v)$ in Example~\ref{ex:lx-n5}, which is
why player $2$ must drop), and exactness supplies it. This suggests the skeleton
\begin{enumerate}[label=(\Alph*),leftmargin=2.3em]
\setcounter{enumi}{3}
\item[(A)] Lemma~\ref{lem:lx-interval}: the bump path is exact;
\item[(D)] each exact $v_\tau$ satisfies $(\star_S)$ at $\phi(v_\tau)$;
\item[(E$'$)] $(\star_S)$ along the whole path implies Lemma~\ref{lem:lx-C},
\end{enumerate}
whose composition is the theorem. The supporting evidence is summarized in
Table~\ref{tab:lx-evidence}; (D) and (E$'$) remain open, and $(\star_S)$ alone is
\emph{not} sufficient (the last row), so a proof of (E$'$) must use exactness along
the whole path, not merely $(\star_S)$ at one endpoint.

\begin{table}[t]
\centering
\small
\caption{Computational evidence (float LP, exact-certifiable) for the reduction.
$(\star_S)$ evaluated at $\phi$ of the base game; ``binding'' bumps only.}
\label{tab:lx-evidence}
\begin{tabular}{@{}p{0.52\linewidth}p{0.22\linewidth}p{0.16\linewidth}@{}}
\toprule
statement & sample & violations\\
\midrule
$\phi$ monotone for sampled proper exact bumps (mixed sizes) & $2000+$ bumps, $n\le6$ & $0$\\
(D) exact $\Rightarrow(\star_S)$ at $\phi$ & $39/39$, $13/13$ & $0$\\
monotone $\Rightarrow(\star_S)$ & all balanced & $0$\\
$(\star_S)$ along the path $\Rightarrow$ monotone & $600$ & $0$\\
$(\star_S)$ at one endpoint $\Rightarrow$ monotone & $1696$ balanced & $22$ non-exact\\
\bottomrule
\end{tabular}
\end{table}

\paragraph{The open problem, sharply.} A Farkas dual turns (D) into a statement
about a single maximal face of the core. For a core point $y$ written as a mixture
of core points, $y(T)\ge v(T)=\phi(T)$ at every tight $T$, so $e=y-\phi$ meets the
tight-cone part of $(\star_S)$ automatically; hence
\begin{align*}
  (\star_S)\ &\Longleftrightarrow\ \exists\,y\in\core(v):\ y_i\ge\phi_i\ (i\in S)
   \ \text{and}\ y(S)>\phi(S)\\
  &\Longleftrightarrow\ \max\{y(S):y\in\core(v),\,y_i\ge\phi_i\ (i\in S)\}>\phi(S).
\end{align*}
Dualizing (verified in $200/200$ random balanced cases, exact-certifiable):
\begin{quote}
$(\star_S)$ \emph{fails} at $\phi(v)$ if and only if $\phi(v)$ maximizes
\[
  \sum_{k\in S}r_k\,x_k
\]
over $\core(v)$ for some weights $r_k\ge1$ on $S$ and zero weights outside
$S$.
\end{quote}
Equivalently, the tight coalitions $\{T:\phi(T)=v(T)\}$ carry nonnegative weights and
a $\beta\ge1$ covering each $k\notin S$ exactly $\beta$ and each $k\in S$ at most
$\beta-1$. Thus \textbf{Lemma~\ref{lem:lx-C}'s hypothesis (D) is exactly the assertion
that an exact game's leximin point never lies on a positive-$S$-weighted maximal face
of its core at a binding-bumpable $S$.} Uniform weights $r_k\equiv c$ are already
excluded, since maximizing $c\,x(S)$ would force $\phi(S)=\max_{\core}x(S)$, contrary
to binding-bumpability; \emph{the entire open content is that exactness excludes the
non-uniform such faces.} The non-exact counterexamples of
Section~\ref{sec:lx-necessity} realize exactly such faces---this is why $(\star_S)$
fails there and a member of $S$ drops. (A corner scan for a deficient $S$-member
pinned at its core-maximum found no occurrence in any exact binding case, $0/16^+$.)

\paragraph{The natural certificate is provably wrong.} Where $(\star_S)$
\emph{does} fail, one would like to prove $d^*_v(S)=0$ and thereby certify
directly that $S$ is not exactly bumpable. This target is
false. From the failure certificate, take the cover
$\{\lambda_U\}$ dual to the failure, $W=\bigcup\{U:\lambda_U>0\}$, and propose
$C=W\cap S$. This route is closed by a six-player exact game (supplement,
\texttt{WITNESS\_KEY\_FALSE\_n6}): $(\star_S)$ fails at $S=\{1,3,4,5\}$ with cover
$\{2,3\},\{1,2,6\},\{3,4,6\}$, yet a core point $x$ with $x(W\cap S)=v(W\cap S)$
attains $x(S)=24>\phi(S)=71/3$, so $W\cap S$ certifies nothing; the associated
slack inequality $\phi(W)-v(W)\le\phi(W\cap S)-v(W\cap S)$ fails by $1/3$. Such
failures are rare ($\approx0.1\%$ of $(\star_S)$-failures at $n\in\{6,7\}$, found
only under generators biased toward thin cores with nested tight chains), which is
why sampling first suggested the cover certificate was universal. In every
observed failure, however, the operative bound survives: there is a coalition $C$
with $M^{(C)}:=\max\{x(S):x\in\core(v),\,x(C)=v(C)\}\le\phi(S)$, so that $S$
cannot be bumped past $\phi(S)$ while preserving exactness (in a $\approx0.2\%$
minority $d^*_v(S)>0$ --- exact six-player witnesses with $d^*=1$ and $d^*=2$ are
included in the supplement --- but always $d^*_v(S)\le\phi(S)-v(S)$). Empirically the
certificate has sharp structure: $C$ can always be taken as a subset of $S$
contained in a \emph{single stratum} of the leximin computation --- almost always
a singleton.  If $C$ is $\phi$-tight, then $\phi\in F_C$, so a certificate can
only be boundary-sharp, $M^{(C)}=\phi(S)$; exact boundary cases of this kind do
occur, while every strict certificate is necessarily non-tight. The open problem is
thus sharpened: produce, from the structure of a $(\star_S)$-failure at the
leximin point of an exact game, a single-stratum $C\subseteq S$ with
$M^{(C)}\le\phi(S)$. Empirically the certificate is pinned further: writing $r$
for the residuals of the failure's dual cover ($r_i\ge1$ on $S$), a
\emph{minimum-residual} member $k\in S$ (one with $r_k=1$) certifies in
$99.6\%$ ($n=6$) to $100\%$ ($n=7$) of observed failures, the residue requiring
a same-stratum pair. By strong duality the singleton form is equivalent to
exhibiting nonnegative weights on $\phi$-tight coalitions with coverage $\beta$
on $N\setminus S$, $\beta-1$ on $S\setminus\{k\}$, and at least $\beta-1$ on
$k$; weights supported inside $S$ suffice in only about two thirds of cases, so
a proof must construct this dual cover from $\phi$-tight coalitions
\emph{crossing} $S$.  The residual normal gives a possible stronger route, but
not an equivalent reformulation.  For a minimum-residual $k$, normalized by
$r_k=1$, and any core $x$ with $x_k=v(\{k\})$, maximality of $\phi$ gives
\[
 [x(S)-\phi(S)]
 +\sum_{i\in S\setminus\{k\}}(r_i-1)(x_i-\phi_i)\leq0.
\]
Consequently the singleton face certificate would follow from the sufficient
inequality
\[
 \sum_{i\in S\setminus\{k\}}(r_i-1)(x_i-\phi_i)\geq0
\]
throughout that face.  The sign is essential; the opposite inequality gives
no upper bound on $x(S)-\phi(S)$.  Whether this stronger face inequality is
forced by the leximin strata is a separate question from the carrier route.
An exact audit answers that stronger question negatively.  A highest-gain
minimum-residual player satisfies the nonnegative-correction condition in
$139/141$ saved multilevel failures, but not in two seven-player cases.  In
one of them
(\texttt{WITNESS\_KEY\_CX\_n7\_s303\_1491}) the correction is identically
$-3/2$ on each of the two highest-gain singleton faces, although both faces
certify and both carrier inequalities hold.  Thus the residual normal alone
cannot prove the multilevel branch; a proof must use additional exactness
geometry such as the off-side carrier comparison.

A new proof-producing exact audit isolates a still sharper target on the saved
hard frontier.  It covers all $141$ stored three-or-more-level failures
($106$ at $n=6$ and $35$ at $n=7$). In every case some minimum-residual
singleton certifies, and, more strongly, some minimum-residual $k$ satisfies
the elementary carrier inequality
\[
 [\phi(P\cup\{k\})-v(P\cup\{k\})]
 \leq [\phi(\{k\})-v(\{k\})].
\]
Across the $367$ individual minimum-residual candidates, $282$ singleton
faces certify and the carrier inequality explains $280$ of them; it never
produces a false certificate.  The remaining two seven-player singleton
certificates require a richer face dual, but another minimum-residual player
in the same case has the carrier certificate.  This is exact finite evidence,
not a theorem.  Nor can it be promoted to an unconditional singleton-carrier
conjecture: the exact two-level witness above has no certifying singleton; its
minimum-residual pair is essential.  At this stage the apparent frontier was a dichotomy.
A failure dual with at most two residual levels is already handled by
Proposition~\ref{prop:lx-two-level-capacity}.  On the genuinely
three-or-more-level branch, these data first suggested proving the displayed
singleton inequality.  The whole-class audit below supersedes that target with
an existential full minimum class $K$ and removes the singleton-selection
problem.  Proposition~\ref{prop:lx-cover-selection} below will show that the
failure cover, and hence $K$, still has to be selected.
The saved corpus initially suggested a canonical gain selector: in all $141$
cases a carrier-certifying $k$ could be chosen to maximize
$\phi_i-v(\{i\})$ inside the minimum-residual class.  That refinement is
false.  A seeded mutation audit produced $534$ further exactly certified
multilevel records ($505$ distinct game--coalition pairs).  Some
minimum-residual carrier still exists in all $534$ records, but a
highest-gain carrier exists in only $520$, and a highest-gain
singleton face certifies in only $533$.

Here is the sharp selector counterexample.  Let the exact six-player game be
the lower envelope of the four common-sum points
\[
\begin{aligned}
 p&=(-6,4,18,-4,11,-10),&q&=(-1,4,24,-4,0,-10),\\
 a&=(-1,-2,19,-4,6,-5),&b&=(-6,0,20,0,9,-10),
\end{aligned}
\]
and take $S=\{1,2,3,4,6\}$.  Its leximin point is
$\phi=(-3,0,20,-2,6,-8)$, the normalized failure cover has three residual
levels and minimum-residual class $K=\{1,2,3\}$, and the singleton gains on
$K$ are $(3,2,2)$.  Thus player $1$ is the unique highest-gain candidate.
Nevertheless its carrier slack is $4>3$, and more strongly
\[
 \max\{x(S):x\in\core(v),\ x_1=v(\{1\})\}=8>7=\phi(S),
\]
so even its singleton face fails.  Players $2$ and $3$ both satisfy the
carrier inequality with equality and certify.  Tight-hypergraph centrality
and cover-weighted adjacency refinements also fail on exact mutations.  Any
singleton route must therefore remain existential: select \emph{some}
minimum-residual player, without imposing a gain ranking.  After adding fresh
three-, four-, and five-generator screens and deduplicating all corpora, this
existential carrier statement holds for all $669$ distinct exact multilevel
game--coalition pairs ($702$ certified records).  The whole-class carrier
statement below supersedes the need to solve this singleton-selection problem.

Proposition~\ref{prop:lx-carrier-separation} sharpens what remains to be
proved.  With $K=\{k\in S:r_k=1\}$, a counterexample would produce
simultaneously the positive-$S$ failure normal $r$ and carrier failure at
every member of $K$.  Proposition~\ref{prop:lx-residual-gap-separator} shows
that the resulting common separator is not a second obstruction to eliminate:
it is automatically the normalized residual gap on $S\setminus K$.  Thus
mere coexistence of the two normals cannot yield a contradiction.  A tempting
coordinatewise strengthening is also false: exactness and the leximin strata
need not permit a convex combination of the carrier-tight faces whose average
displacement is nonnegative at every higher-residual coordinate.  In the exact
six-player instance \texttt{WITNESS\_KEY\_CX\_n6\_s201\_10919}, the best common
coordinatewise margin over all such mixtures is $-1/20$, although players $4$
and $5$ both satisfy the carrier inequality (each with margin $-3/4$).

The scalar condition in Corollary~\ref{cor:lx-residual-gap-escape} is the
minimal sufficient statement for a \emph{fixed} failure cover, and it initially
survives all $141$ saved hard multilevel failures.  Across their $367$
minimum-residual faces, $277$ escape; all $87$ failed carriers are strictly
negative, while three valid carriers do not escape but another candidate does
in the same record.

Adversarial mutation shows that even this fixed-cover target is not universal.
Among $505$ distinct exact mutation cases, all have a minimum-residual carrier
but only $123$ have an escape face for the cover returned by the exact failure
dual.  Already trial $0$ of the first seeded face-selector mutation audit has
$K=\{1,2,3\}$ and normalized residual-gap weights $(2/3,1/3)$ on players
$5,6$.  Players $1$ and $2$ satisfy the carrier inequality, with margins
$-2$ and $0$, yet their face maxima are both $-1/3$; player $3$ fails the
carrier inequality and has maximum $-4/3$.  Consequently neither
coordinatewise neutralization nor escape for an arbitrarily chosen failure
cover can be the missing theorem.  A residual-based proof would have to select
a favorable cover from the full dual face jointly with the carrier candidate;
otherwise it must return to a cover-independent exactness exchange.

There is one dimension in which this cover-selection problem disappears.

\begin{lemma}[three-player residual-order rigidity]
\label{lem:lx-three-player-order}
Suppose $|S|=3$ and $(\star_S)$ fails at a core point $\psi$.  If no failure
cover has at most two residual levels on $S$, then every failure cover has
three distinct residuals, and all failure covers induce the same strict order
on $S$.  In particular, their minimum-residual class is the same singleton.
\end{lemma}

\begin{proof}
Before the harmless positive normalization, failure covers form a polyhedral
cone; the normalized residual set is convex.  A residual tie on a three-player
set has at most two distinct levels, so every residual vector in the stated
branch has a strict order.  If two covers induced different orders, their line
segment would remain in the failure-cover set.  Along that
segment the order could change only at a residual tie, producing a cover with
at most two levels, a contradiction.
\end{proof}

\begin{lemma}[three-player block-star and residual-gap alternatives]
\label{lem:lx-three-player-block-star}
Let $\mathcal T(\psi)$ be the proper coalitions tight at a core point $\psi$
and put
\[
 \mathcal D_\psi=\{d\in\mathbb R^N:d(N)=0,\ d(T)\geq0
       \text{ for every }T\in\mathcal T(\psi)\}.
\]
For a pair $\{a,b\}\subset S$ with third player $c$, no failure cover has
$r_a=r_b$ if and only if some $d\in\mathcal D_\psi$ satisfies
\[
 d_c\geq0,\qquad d_a+d_b\geq0,\qquad d(S)>0.       \tag{BS}
 \label{eq:lx-block-star}
\]
Consequently, when $|S|=3$, every failure cover has three distinct residuals
if and only if all three singleton--complementary-pair versions of
\eqref{eq:lx-block-star} hold.

Now suppose failure covers exist and every cover has the fixed order
$r_k<r_i<r_j$.  Normalize by
$r_k=1$ and define
\[
 \delta=\min\{r_j-r_i:r\text{ is a failure-cover residual vector},\ r_k=1\}.
 \tag{AG}\label{eq:lx-adjacent-gap}
\]
Then $\delta>0$, and there is $d\in\mathcal D_\psi$ with
\[
 (d_k,d_i,d_j)=(\delta,1,-1).
 \tag{GD}\label{eq:lx-gap-direction}
\]
In particular, if $P=N\setminus S$ and $A=P\cup\{k\}$, then
$d(P)=-\delta$ and $d(A)=0$.
\end{lemma}

\begin{proof}
For every failure cover and every $d\in\mathcal D_\psi$, the cover identity
gives
\[
 \sum_{p\in S}r_pd_p
 =\mu d(N)-\sum_{T\in\mathcal T(\psi)}\lambda_Td(T)\leq0. \tag{P}
 \label{eq:lx-cover-tangent-pairing}
\]
If $r_a=r_b$, the left side is
$r_a(d_a+d_b)+r_cd_c$, which is strictly positive under
\eqref{eq:lx-block-star}.  This proves one direction.  Conversely, apply
Farkas' lemma to the normalized failure-cover system with the additional
equality $r_a=r_b$.  The player-equation multipliers, with sign reversed,
give $d\in\mathcal D_\psi$; the multipliers of the lower bounds on $r_c$ and
on the tied pair give respectively $d_c\geq0$ and $d_a+d_b\geq0$.  Normalizing
the Farkas right-hand side gives $d(S)>0$.  This proves the alternative.

For the second claim, the polyhedron in \eqref{eq:lx-adjacent-gap} is nonempty.
Its objective is nonnegative and attains its minimum.  If the minimum were
zero, an optimizer would have $r_i=r_j$, contrary to the fixed strict order;
hence $\delta>0$.  Dualize this LP.  The dual player multipliers again define
a vector in $\mathcal D_\psi$.  At every feasible point
$r_j>r_i>r_k=1$, so the lower bounds on $r_i$ and $r_j$ are slack; complementary
slackness and the objective $r_j-r_i$ give $d_i=1$ and $d_j=-1$.
The dual price of the normalization $r_k=1$, together with equality of the
primal and dual values, gives $d_k=\delta$.  Finally
$d(S)=\delta$ and $d(N)=0$, so $d(P)=-\delta$ and $d(A)=0$.
\end{proof}

Thus the unresolved $|S|=3$ branch has no chamber-choice ambiguity: after
Proposition~\ref{prop:lx-two-level-capacity}, it reduces to proving the carrier
inequality for one uniquely determined minimum-residual player.  This is
strictly simpler than the four-player geometry below, where the residual cone
can cross several genuinely multilevel chambers.

Lemma~\ref{lem:lx-three-player-block-star} replaces the universal quantifier
over the cover polyhedron by three primal tangent directions.  It also supplies
one canonical direction parallel to the trace face $A=P\cup\{k\}$.  In the six
exact hard cases, independently verified rational Farkas certificates exist
for all $18$ pair-tie systems.  Moreover every normalized adjacent gap is
$\delta=1$, and the six directions in \eqref{eq:lx-gap-direction} are integral.
These values are finite evidence; the block-star and residual-gap alternatives
themselves are proved for every three-player target.

There is a natural stronger three-player target.  If $k$ is that unique
minimum-residual player, ask whether
\[
 \max\{x(P):x\in\core(v),\ x(P\cup\{k\})=v(P\cup\{k\})\}
 \geq\phi(P).                                      \tag{TF}
 \label{eq:lx-three-player-face}
\]
Any maximizer in \eqref{eq:lx-three-player-face} is an exact-face witness in
Lemma~\ref{lem:lx-generator-compensation}, so both terms on the right side of
\eqref{eq:lx-generator-compensation} are nonnegative and the singleton carrier
inequality follows.  The initial exact chamber audit contained six
genuinely multilevel $|S|=3$ cases.  Each had exactly one accessible
minimum-class chamber, a strict carrier margin between $-8$ and $-1/2$, and
an exact rational primal--dual certificate for \eqref{eq:lx-three-player-face}.
A targeted seven-player screen of $1{,}976$ nondegenerate lower-envelope games
found $649$ binding star failures and seven float multilevel candidates.  Four
extremals were recomputed exactly: one collapsed to the proved two-level
branch, while the other three supplied three of the six cases above.  These
counts are evidence; Lemma~\ref{lem:lx-three-player-order} is the proved
dimension reduction.

Those six face duals share a sharper pattern.  Each has zero grand-coalition
price, and every positively weighted core constraint $U$ satisfies
$U\cap S=\{k\}$.  If $b>0$ is minus the price of the face equality and
$\lambda_U\geq0$ are the resulting core weights, stationarity reads
\[
 \sum_U\lambda_U\mathbf1_U
 =(b-1)\mathbf1_P+b\mathbf1_{\{k\}},
\]
and \eqref{eq:lx-three-player-face} is exactly the weighted slack inequality
\[
 \sum_U\lambda_U[\phi(U)-v(U)]
 \geq b[\phi(P\cup\{k\})-v(P\cup\{k\})].       \tag{TS}
 \label{eq:lx-three-player-slack-dual}
\]
This common dual support pattern was exact finite evidence, not a theorem.
The counterexample below has a positive dual coalition disjoint from $S$ and
therefore also falsifies the exact-trace-$\{k\}$ support refinement.  The next
reduction remains valid, but its proposed sign conclusion will likewise be
falsified.

Exactness supplies a direct face-coordinate reduction.  It neither assumes the
observed trace-face dual nor selects a dual representative.

\begin{lemma}[three-player residual-spread bound]
\label{lem:lx-three-player-spread}
Let $v$ be exact and suppose a normalized failure cover has residuals
$r_k=1<r_i=\rho<r_j=\sigma$ on $S=\{k,i,j\}$, and put
$A=P\cup\{k\}$.  For every $a\in F_A$,
\[
 a(P)-\phi(P)
 \geq(\rho-1)[\phi(A)-v(A)]
      +(\sigma-\rho)(a_j-\phi_j).             \tag{FC}
 \label{eq:lx-three-player-face-coordinate}
\]
Consequently, writing $M_j^A=\max_{a\in F_A}a_j$, the target margin satisfies
\[
 \begin{aligned}
 \max_{x\in F_A}x(P)-\phi(P)
 &\geq (\rho-1)[\phi(A)-v(A)]\\
 &\quad +(\sigma-\rho)(M_j^A-\phi_j)\\
 &\geq (\rho-1)[\phi(A)-v(A)]\\
 &\quad -(\sigma-\rho)[\phi_j-v(\{j\})].
 \end{aligned}                                      \tag{RS}
 \label{eq:lx-three-player-spread}
\]
In particular, \eqref{eq:lx-three-player-face} holds whenever the right side of
\eqref{eq:lx-three-player-spread} is nonnegative.
\end{lemma}

\begin{proof}
Write the tight-coalition failure cover as
\[
 (1,\rho,\sigma)_S
 =\mu\mathbf1_N-\sum_T\alpha_T\mathbf1_T,
 \qquad \alpha_T\geq0,
\]
where $\phi(T)=v(T)$ whenever $\alpha_T>0$.  For $a\in F_A$, put
$\Delta=a-\phi$ and
$s(A)=\phi(A)-v(A)$.  Efficiency and tightness at $A$ give
\[
 \Delta(P)+\Delta_k=-s(A),\qquad
 \Delta_i+\Delta_j=s(A).
\]
Pairing the failure cover with $\Delta$, and using core feasibility of $a$ on
every positively weighted $\phi$-tight coalition, gives
\[
 \Delta_k+\rho\Delta_i+\sigma\Delta_j
 =-\sum_T\alpha_T[a(T)-v(T)]\leq0.
\]
Substituting $\Delta_i=s(A)-\Delta_j$ and
$\Delta(P)=-s(A)-\Delta_k$ proves
\eqref{eq:lx-three-player-face-coordinate}.  Maximize $a_j$ over $F_A$.
Finally every core point has $a_j\geq v(\{j\})$, which gives
\eqref{eq:lx-three-player-spread}.
\end{proof}

The singleton-floor comparison in \eqref{eq:lx-three-player-spread} is strict
in only two of the original six exact three-level cases and has lower bounds
$-5/4,-5/3,-3,-2$ in the other four.  The stronger face-coordinate bound
\eqref{eq:lx-three-player-face-coordinate} proves all six.  Its exact values
are
\[
 \frac12,\quad1,\quad\frac{11}{4},\quad\frac43,\quad0,\quad2,
\]
versus actual target margins
$1/2,1,15/4,4/3,5,2$; hence it is sharp in four cases.  A natural attempted
generalization was the one-coordinate inequality
\[
 M_j^A-\phi_j
 \geq-\frac{\rho-1}{\sigma-\rho}[\phi(A)-v(A)]. \tag{HC}
 \label{eq:lx-high-residual-coordinate}
\]
Unlike trace extension itself, \eqref{eq:lx-high-residual-coordinate} asks only
for the high-residual coordinate on one exact face.  It is nevertheless false.

\begin{proposition}[genuinely multilevel face-target counterexample]
\label{prop:lx-three-player-face-target-counterexample}
Let $v$ be the lower envelope of the three common-sum allocations
\[
 \begin{pmatrix}
 14&-11&12&-16&9&14&10\\
 15&-11&12&-16&8&14&10\\
 18& -6& 6&-17&6&13&12
 \end{pmatrix}.
\]
Then $v$ is an exact seven-player game.  For $S=\{2,3,5\}$ its leximin
point is
\[
 \phi=\left(\frac{47}{3},-9,\frac{29}{3},-\frac{49}{3},
              \frac{23}{3},\frac{41}{3},\frac{32}{3}\right).
\]
The coalition is binding by $5/3$, $(\star_S)$ fails, and every failure cover
has three distinct residuals in the order
\[
 r_5<r_3<r_2.
\]
Nevertheless, for $P=N\setminus S$ and $A=P\cup\{5\}$,
\[
 \max_{x\in F_A}x(P)=23<\frac{71}{3}=\phi(P).       \tag{FCX}
 \label{eq:lx-face-target-counterexample}
\]
For the normalized cover $(r_5,r_3,r_2)=(1,2,3)$ one has
$s(A)=1/3$ and $M_2^A-\phi_2=-1$, so both
\eqref{eq:lx-three-player-face} and
\eqref{eq:lx-high-residual-coordinate} fail by $2/3$.
Trace extension itself still holds with gap zero, and the desired carrier
inequality still holds with margin
$s(A)-s(\{5\})=-4/3$.
\end{proposition}

\begin{proof}
All three rows have sum $32$ and are core points of their lower envelope.
Every proper-coalition worth is attained by a row, proving exactness.  Exact
leximin LPs give successive surplus levels
$2/3,5/3,2,11/3$, pinning respectively
$\{4,6,7\},\{1,5\},\{2\},\{3\}$.  A core LP gives the stated binding
margin.  The tight-cover weights $1,1,1,2$ on coalition masks
$25,49,69,104$, with grand coefficient $3$, give residuals
$(r_2,r_3,r_5)=(3,2,1)$ and certify star failure.  Exact Farkas certificates
make each of the three pair-tie cover systems infeasible, so the whole cover
face is genuinely multilevel by Lemma~\ref{lem:lx-three-player-order}.

On $F_A$, the point
$(15,-10,11,-16,8,14,10)$ attains both $x(P)=23$ and $x_2=-10$.
For the $P$-maximum, unit core-dual weights on masks $25,49,80,104$, grand
price zero, and face price $-3$ certify $23$.  A coordinate dual certifies
$M_2^A=-10$: its core weights on masks $25,49,69,80,104$ are respectively
$2,2,1,1,3$, with grand price $-1$ and face price $-4$.  Direct substitution
gives the two margins in the statement.
The restricted $A$-core has the same $P$-maximum $23$, proving zero trace
gap, while the singleton and carrier slacks give $-4/3$.
\end{proof}

The failure is not a numerical knife edge.  Replace the second generator by
\[
 p^2(t)=(17-t,-11,12,-16,6+t,14,10).
\]
Exact recomputation at
\[
 t=0,\frac12,1,\frac43,\frac32,2,3
\]
keeps the leximin point, binding margin $5/3$, star opening zero, strict
three-level residual order, carrier margin $-4/3$, and trace gap zero fixed.
The face-target margins are respectively
\[
 \frac43,\quad\frac56,\quad\frac13,\quad0,
 \quad-\frac16,\quad-\frac23,\quad-\frac23.
\]
Thus an exact sum-preserving motion within $A$ crosses the stronger target
without moving the carrier comparison at all.  These rational samples are a
path audit, not a claim about every real parameter value.

The failed one-coordinate target has an exact dual form that separates the
face price from coalition-slack compensation and diagnoses the crossing.

\begin{proposition}[high-coordinate slack--price identity]
\label{prop:lx-high-coordinate-dual}
Let $F_A$ be nonempty, put $s(C)=\phi(C)-v(C)$, and let
$M_j^A=\max_{a\in F_A}a_j$.  Every optimal dual for this coordinate maximum
can be written with $\lambda_C\geq0$ and $u,b\in\mathbb R$ as
\[
 -\mathbf1_{\{j\}}=
 \sum_{\emptyset\ne C\subsetneq N}\lambda_C\mathbf1_C
       +u\mathbf1_N+b\mathbf1_A,                    \tag{CD}
 \label{eq:lx-coordinate-dual}
\]
and satisfies the exact identity
\[
 M_j^A-\phi_j=\sum_C\lambda_Cs(C)+b\,s(A).          \tag{SP}
 \label{eq:lx-coordinate-slack-price}
\]
Consequently, for any $c\geq0$,
\[
 M_j^A-\phi_j\geq-cs(A)
 \quad\Longleftrightarrow\quad
 \sum_C\lambda_Cs(C)+(b+c)s(A)\geq0.               \tag{DP}
 \label{eq:lx-coordinate-dual-price}
\]
The equivalence holds for one, hence for every, optimal coordinate dual.  In
particular $b\geq-c$ is sufficient; if $b<-c$, only the positively weighted
coalition slacks can supply the missing amount.
\end{proposition}

\begin{proof}
Write the core inequalities as $-a(C)\leq-v(C)$ and minimize $-a_j$ subject
also to grand efficiency and $a(A)=v(A)$.  Reversing the signs of the
nonpositive inequality-dual variables gives \eqref{eq:lx-coordinate-dual},
while strong duality gives
\[
 -M_j^A=\sum_C\lambda_Cv(C)+u\,v(N)+b\,v(A).
\]
Pair \eqref{eq:lx-coordinate-dual} with $\phi$, use
$\phi(N)=v(N)$, and subtract the dual objective identity.  This proves
\eqref{eq:lx-coordinate-slack-price}; adding $cs(A)$ proves
\eqref{eq:lx-coordinate-dual-price}.  Its left side is independent of the
chosen optimal dual, which proves the final assertion.
\end{proof}

In the original six exact hard faces, maximizing the face price over the optimal dual
face gives
\[
 -1,\quad-\frac23,\quad0,\quad-2,\quad0,\quad0.
\]
Here $c=(\rho-1)/(\sigma-\rho)=1$.  Thus the price condition alone proves five
cases.  In the fourth, the maximal price is $-2$, but the weighted slack is
$5/3$ and $s(A)=1/3$, so the right side of
\eqref{eq:lx-coordinate-dual-price} is $4/3>0$.  In
Proposition~\ref{prop:lx-three-player-face-target-counterexample}, the maximal
price is $-4$, the weighted slack is $1/3$, and $s(A)=1/3$; hence the same
right side is $-2/3$.  Thus the identity distinguishes genuine compensation
from a genuine sign failure, while a face-price bound by itself was already
false on the positive original corpus.

The closed-form spread bound is one feasible point of a stronger finite
exactness optimization.  Write $s(C)=\phi(C)-v(C)$ and, for a trace-face dual
$(\lambda,b)$, define
\[
 \begin{split}
 \mathcal B(\lambda,A)=\sup\{&(c-b)s(A)-\sum_C\eta_Cs(C):
       \eta_C\geq0,\ c\geq b,\ q\in\mathbb R,\\
   &\sum_U\lambda_U\mathbf1_U+
       \sum_C\eta_C\mathbf1_C-c\mathbf1_A=q\mathbf1_N\}.
 \end{split}
 \tag{OB}\label{eq:lx-optimal-exact-balance}
\]

\begin{proposition}[one-face exact-balance bound]
\label{prop:lx-optimal-exact-balance}
Under the hypotheses of Lemma~\ref{lem:lx-three-player-spread}, suppose in
addition that the face LP has a zero-grand dual $(\lambda,b)$ satisfying
\[
 \sum_U\lambda_U\mathbf1_U
 =(b-1)\mathbf1_P+b\mathbf1_{\{k\}}.
\]
Then
\[
 \max_{x\in F_A}x(P)-\phi(P)\geq\mathcal B(\lambda,A).
\]
The singleton-floor residual-spread right side in
\eqref{eq:lx-three-player-spread} is obtained by a feasible choice in
\eqref{eq:lx-optimal-exact-balance}.
\end{proposition}

\begin{proof}
Fix a feasible $(\eta,c,q)$ and choose, by exactness,
$a\in\core(v)$ with $a(A)=v(A)$.  Dot the incidence identity in
\eqref{eq:lx-optimal-exact-balance} first with $a$ and then with $\phi$.
Core feasibility for every positively weighted coalition and tightness at
$A$ give
\[
 \sum_U\lambda_Us(U)+\sum_C\eta_Cs(C)-cs(A)\geq0.
\]
Subtracting $bs(A)$ instead of $cs(A)$ proves the asserted bound and taking
the supremum completes the first claim.  For the second, use the tight
failure-cover weights as part of $\eta$, add
$(\sigma-\rho)$ on the singleton $\{j\}$, and take
$c=b+\rho-1$ and $q=\mu-\rho$.  Feasibility is the incidence identity
\[
 \sum_U\lambda_U\mathbf1_U+
 \sum_T\alpha_T\mathbf1_T+(\sigma-\rho)\mathbf1_{\{j\}}
 -(b+\rho-1)\mathbf1_A=(\mu-\rho)\mathbf1_N.
\]
\end{proof}

The exact LP for \eqref{eq:lx-optimal-exact-balance} gives, in the six known
three-level cases, respectively
\[
 \frac12,\quad 1,\quad \frac{15}{4},\quad-\frac53,\quad-2,\quad 2.
\]
It therefore proves the full face target in four cases and equals the actual
target margin in all four, including the case where the closed-form spread
bound is $-5/4$ but the target margin is $15/4$.  Restricting the auxiliary
weights $\eta$ to $\phi$-tight coalitions is infeasible in all six cases, so
non-tight slack terms are indispensable.  The last two actual target margins
exceed their optimal one-face bounds by $3$ and $7$, but neither remains an
unexplained finite case.  In the latter, $s(A)=0$, so $\phi\in F_A$ proves the
face target immediately.  In the former, a common-sum generator attaining
$v(A)$ has outside sum $25>71/3=\phi(P)$ and is the required exact-face
witness.  Thus all six cases now have proof-producing decompositions; only
the last generator argument was specific to the finite lower-envelope model.
The unconditional face-coordinate bound above now supplies exact-face proofs
for both remaining cases as well, so none of the six finite proofs requires a
model-specific generator argument.

The remaining dual-existence question has a clean primal form.  For
$A=P\cup\{k\}$ write
\[
 \core_A(v)=\{y\in\mathbb R^A:y(A)=v(A),\ y(B)\geq v(B)
                    \text{ for every }B\subsetneq A\}.
\]
The projection of $F_A$ is contained in $\core_A(v)$, so always
\[
 \max_{x\in F_A}x(P)\leq\max_{y\in\core_A(v)}y(P).       \tag{TE}
 \label{eq:lx-trace-extension}
\]
Equality is precisely the existence of an optimal restricted allocation that
extends to the full core.  The corresponding dual statement is exact.

\begin{lemma}[trace-extension dual equivalence]
\label{lem:lx-trace-extension-dual}
Equality holds in \eqref{eq:lx-trace-extension} if and only if the full face
LP has an optimal dual with zero grand-coalition price.  Every such dual can
be chosen with all positively weighted core constraints contained in $A$ and
has
\[
 \sum_U\lambda_U\mathbf1_U
 =(b-1)\mathbf1_P+b\mathbf1_{\{k\}}
\]
for some $b\geq1$.
\end{lemma}

\begin{proof}
The restricted problem retains only the coordinates and core constraints
inside $A$, together with efficiency on $A$.  Any restricted dual lifts to
the full face dual by assigning zero price to grand efficiency and zero
weight to every other core constraint.  Conversely, stationarity for a
zero-grand full dual has zero coverage on $S\setminus\{k\}$.  Since core
weights are nonnegative, every positively weighted coalition is contained in
$A$; deleting the unused coordinates gives a restricted dual of the same
value.  Strong duality proves the equivalence.  Stationarity for the
$P$-objective and the face equality gives the displayed balance, whose
nonnegative coverage forces $b\geq1$.
\end{proof}

Thus trace extension supplies the balance identity used in
Proposition~\ref{prop:lx-optimal-exact-balance}.  The unconditional
Lemma~\ref{lem:lx-three-player-spread} no longer requires trace extension or
any choice of face dual; the stronger observed fact that every positive dual
coalition has trace exactly $\{k\}$ is not required there either.

For the three-player branch the primal obstruction collapses to one scalar
inequality.  Put $R=N\setminus A=\{i,j\}$ and
$d_A=v(N)-v(A)$.  For $y\in\core_A(v)$ and nonempty $D\subseteq R$, define
the outside demand
\[
 \ell_D(y)=\max_{B\subseteq A}\{v(B\cup D)-y(B)\}.
\]

\begin{proposition}[two-player-complement trace criterion]
\label{prop:lx-two-player-trace}
Let $v$ be exact, let $A=N\setminus\{i,j\}$, and let
$y\in\core_A(v)$.  Then $y$ extends to a point of $F_A$ if and only if
\[
 \ell_{\{i\}}(y)+\ell_{\{j\}}(y)\leq d_A.       \tag{SD}
 \label{eq:lx-split-demand}
\]
In particular, exactness rules out every obstruction involving a coalition
that contains both omitted players.
\end{proposition}

\begin{proof}
An extension chooses two coordinates $z_i,z_j$ with $z_i+z_j=d_A$.
Constraints whose traces on $R$ are $\{i\}$ or $\{j\}$ are exactly
$z_i\geq\ell_{\{i\}}(y)$ and $z_j\geq\ell_{\{j\}}(y)$, so such coordinates
exist exactly when \eqref{eq:lx-split-demand} holds.  Constraints with empty
trace already hold because $y\in\core_A(v)$.

It remains to check a trace equal to $R$.  For each $B\subseteq A$, exactness
provides $a\in\core(v)$ with $a(B)=v(B)$.  Core feasibility and efficiency
give
\[
 v(B\cup R)-v(B)\leq a(R)=v(N)-a(A)\leq v(N)-v(A)=d_A.
\]
Since $y(B)\geq v(B)$, this implies
$v(B\cup R)-y(B)\leq d_A$.  Thus every two-player-trace constraint is
automatic, completing the equivalence.
\end{proof}

Consequently equality in \eqref{eq:lx-trace-extension} is equivalent to the
restricted face LP having an optimum $y$ that satisfies (SD).  This replaces
a high-dimensional extension question by the exclusion of two witness
coalitions, one carrying $i$ and one carrying $j$.

There is an exact dual for that last selection problem.  Put
\[
 m_A=\min_{y\in\core_A(v)}y_k,\qquad
 \Gamma_A=\min_{\substack{y\in\core_A(v)\\y_k=m_A}}
 \{\ell_{\{i\}}(y)+\ell_{\{j\}}(y)-d_A\}.
\]
Then equality in \eqref{eq:lx-trace-extension} is equivalent to
$\Gamma_A\leq0$.

\begin{lemma}[split-obstruction dual]
\label{lem:lx-split-obstruction-dual}
The number $\Gamma_A$ is the maximum of
\[
 \sum_C\tau_Cv(C)+
 \sum_{p\in\{i,j\}}\sum_{B\subseteq A}\theta_{pB}v(B\cup\{p\})
 -v(N)-(\alpha-1)v(A)-\beta m_A                         \tag{SDD}
 \label{eq:lx-split-obstruction-dual}
\]
over $\tau_C,\theta_{pB}\geq0$ and $\alpha,\beta\in\mathbb R$ satisfying
\[
 \sum_B\theta_{pB}=1\quad(p=i,j)
\]
and the lifted incidence balance
\[
 \sum_C\tau_C\mathbf1_C+
 \sum_{p\in\{i,j\}}\sum_{B\subseteq A}
       \theta_{pB}\mathbf1_{B\cup\{p\}}
 =\mathbf1_N+(\alpha-1)\mathbf1_A+\beta\mathbf1_{\{k\}}.
 \tag{SDB}\label{eq:lx-split-obstruction-balance}
\]
Here $C$ ranges over the proper coalitions contained in $A$.
\end{lemma}

\begin{proof}
Minimize $u_i+u_j-d_A$ subject to $y\in\core_A(v)$, $y_k=m_A$, and
\[
 y(B)+u_p\geq v(B\cup\{p\})
 \qquad(B\subseteq A,\ p=i,j).
\]
At optimum $u_p=\ell_{\{p\}}(y)$.  In the LP dual, stationarity of $u_p$
gives $\sum_B\theta_{pB}=1$.  Stationarity of the coordinates of $y$ first
gives the corresponding balance on the traces $B$; replacing each trace by
$B\cup\{p\}$ adds one unit on each omitted player and yields
\eqref{eq:lx-split-obstruction-balance}.  The dual objective, after
subtracting $d_A=v(N)-v(A)$, is
\eqref{eq:lx-split-obstruction-dual}.  Strong duality completes the proof.
\end{proof}

The split dual can be compared with the simpler dual for $m_A$.  Write that
dual as
\[
 \sum_C\kappa_C\mathbf1_C=c\mathbf1_A+\mathbf1_{\{k\}},
 \qquad
 \sum_C\kappa_Cv(C)=cv(A)+m_A,                         \tag{MD}
 \label{eq:lx-minimum-coordinate-dual}
\]
with $\kappa_C\geq0$ on proper coalitions contained in $A$.

The coefficient condition also has a one-dimensional primal form.  Fix a
minimum-coordinate dual coefficient $c$ and define the extended-real value
\[
 \begin{split}
 \widehat\Gamma_{A,c}(t)=\min\{&u_i+u_j+t-d_A:\
       y(A)=v(A)+t,\ y_k=m_A-ct,\\
   &y(B)\geq v(B)\quad(B\subsetneq A),\qquad
       y(B)+u_p\geq v(B\cup\{p\})
       \quad(B\subseteq A,\ p=i,j)\}.
 \end{split}
 \tag{PT}\label{eq:lx-minimum-dual-tangent}
\]
and put the value equal to $+\infty$ when the constraints are infeasible.
Thus $\widehat\Gamma_{A,c}(0)=\Gamma_A$.

\begin{proposition}[minimum-dual tangent criterion]
\label{prop:lx-minimum-dual-tangent}
The optimal split-dual face contains a member satisfying
$\alpha-1=\beta c$ if and only if
\[
 \widehat\Gamma_{A,c}(t)\geq\Gamma_A
 \qquad\text{for every }t\in\mathbb R.       \tag{TN}
 \label{eq:lx-minimum-dual-tangent-nondescent}
\]
Equivalently, compatibility says that $t=0$ globally minimizes the convex
parametric split value in \eqref{eq:lx-minimum-dual-tangent}.
\end{proposition}

\begin{proof}
Write $\gamma=\alpha-1$ and let $J$ be the split-dual objective at $t=0$.
The dual feasible region is independent of $t$, while its objective at $t$ is
\[
 J+t(\beta c-\gamma).
\]
Hence $\widehat\Gamma_{A,c}$ is the supremum of affine functions and is
convex.  Its subdifferential at zero is the set of slopes
$\beta c-\gamma$ attained on the optimal dual face: that face is convex and
the slope map is affine, so its image is already an interval.  Therefore
$0\in\partial\widehat\Gamma_{A,c}(0)$ exactly when some optimal dual has
$\gamma=\beta c$.  The convex optimality criterion identifies this with
\eqref{eq:lx-minimum-dual-tangent-nondescent}.
\end{proof}

\begin{proposition}[compatible minimum-cover certificate]
\label{prop:lx-compatible-minimum-cover}
Suppose an optimal split dual in
Lemma~\ref{lem:lx-split-obstruction-dual} and an optimal minimum-coordinate
dual in \eqref{eq:lx-minimum-coordinate-dual} satisfy
\[
 \alpha-1=\beta c.
\]
If there is $z\in\core(v)$ with $z(C)=v(C)$ whenever $\kappa_C>0$, then
$\Gamma_A\leq0$ and trace extension holds.  When $v$ is exact and $c=0$, the
common-tightness condition is automatic.
\end{proposition}

\begin{proof}
Let $\omega$ denote the aggregate coalition weights on the left side of
\eqref{eq:lx-split-obstruction-balance}.  The coefficient identity and (MD)
give the signed incidence identity
\[
 \sum_C\omega_C\mathbf1_C-
 \beta\sum_C\kappa_C\mathbf1_C=\mathbf1_N.
\]
Because $\omega_C\geq0$, core feasibility can only increase the $\omega$-sum
when worths are replaced by $z$-values.  Common tightness leaves the
$\kappa$-sum unchanged, irrespective of the sign of $\beta$.  Hence
\[
 \sum_C\omega_Cv(C)-\beta\sum_C\kappa_Cv(C)
 \leq
 \sum_C\omega_Cz(C)-\beta\sum_C\kappa_Cz(C)=v(N).
\]
Using (MD) and $\alpha-1=\beta c$, the left side minus $v(N)$ is exactly the
optimal objective \eqref{eq:lx-split-obstruction-dual}, namely $\Gamma_A$.
Thus $\Gamma_A\leq0$.

If $c=0$, the nonnegative incidence identity in (MD) forces
$\kappa_{\{k\}}=1$ and every other $\kappa_C=0$.  Exactness at $\{k\}$ then
supplies the required $z$.
\end{proof}

The other hypothesis also has a scalar primal form.  For a fixed optimal
minimum cover $\kappa$, its positive support is jointly tight at a full-core
point if and only if
\[
 \min_{z\in\core(v)}\{c z(A)+z_k\}=c v(A)+m_A.       \tag{NL}
 \label{eq:lx-minimum-normal-lift}
\]
Indeed, (MD) makes the left side at least the right side, and equality is
equivalent to zero total $\kappa$-weighted core slack.  Thus the remaining
compatible-cover route is entirely primal: prove tangent non-descent
\eqref{eq:lx-minimum-dual-tangent-nondescent} and the minimum-normal lift
\eqref{eq:lx-minimum-normal-lift} from star failure and the leximin strata.
These two statements concern the same face.  Every feasible $y$ in
\eqref{eq:lx-minimum-dual-tangent} satisfies
\[
 \sum_C\kappa_C[y(C)-v(C)]
 =c[y(A)-v(A)]+(y_k-m_A)=ct-ct=0.
 \tag{MF}\label{eq:lx-minimum-cover-face}
\]
All summands are nonnegative, so every positive $\kappa$-coalition is tight
along the entire tangent LP.  Hence tangent non-descent is an optimization
statement on the restricted minimum-cover face, while
\eqref{eq:lx-minimum-normal-lift} says that this same tight face meets the full
core.

This also gives an entirely primal proof of
Proposition~\ref{prop:lx-compatible-minimum-cover}.  Let $z$ attain
\eqref{eq:lx-minimum-normal-lift} and put $t=z(A)-v(A)$.  Equality in (NL)
gives $z_k=m_A-ct$, so the restriction of $z$ to $A$, together with
$u_i=z_i$ and $u_j=z_j$, is feasible in (PT).  Full efficiency gives
$z_i+z_j=d_A-t$, hence this feasible point has objective zero.  Therefore
\[
 \Gamma_A=\widehat\Gamma_{A,c}(0)
 \leq\widehat\Gamma_{A,c}(t)\leq0
\]
by tangent non-descent.

Exactness alone does not force the normal lift, even at unit price.  Define
\[
 \Delta^{\rm NL}_{A,k}
 =\min_{z\in\core(v)}\{z(A)+z_k\}-[v(A)+m_A]
 \qquad(c=1).
\]

The most natural quantitative comparison would be
\[
 \Delta^{\rm NL}_{A,k}\leq\eta_S(\phi).
 \tag{SN}\label{eq:lx-star-to-normal}
\]
It is false, and its sharp-bound deformation also refutes star-conditioned
trace extension itself.

\begin{proposition}[exact star--trace counterexample]
\label{prop:lx-star-trace-counterexample}
For $0\leq\varepsilon\leq1$, put
\begin{align*}
 p&=(13,13,-12,-3,13,0,1),\\
 q^\varepsilon&=(11,14+\varepsilon,-3,2-\varepsilon,16,-10,-5),\\
 u&=(19,19,-20,-3,20,2,-12),
\end{align*}
and let $v^\varepsilon$ be their lower-envelope game, with common grand worth
$25$.  Let $S=\{2,4,5\}$, $k=5$, $A=\{1,3,5,6,7\}$, and
$P=A\setminus\{5\}$.

At $\varepsilon=1/10$, the restricted minimum has a unit-price dual but
\[
 \Delta^{\rm NL}_{A,5}=1>\frac{17}{20}=\eta_S(\phi),
\]
so \eqref{eq:lx-star-to-normal} fails.  At $\varepsilon=2/3$, the star
condition fails at the leximin point, $S$ is binding, and nevertheless
\[
 \max_{y\in\core_A(v)}y(P)=-5>-6=
 \max_{x\in F_A}x(P).
 \tag{CX}\label{eq:lx-star-trace-counterexample}
\]
Thus trace extension fails, with trace gap $1$ and split obstruction
$\Gamma_A=2$.
\end{proposition}

\begin{proof}
All three generators have sum $25$ and belong to the core of their
lower-envelope game.  Every proper-coalition worth is attained by a generator,
so $v^\varepsilon$ is exact.

At $\varepsilon=1/10$, exact primal--dual LP pairs give restricted minimum
$14$, the unit cover
$\kappa_{\{3,5,6\}}=\kappa_{\{1,5,7\}}=1$, normal-lift gap $1$, and star
opening $17/20$.

At $\varepsilon=2/3$, the leximin point is
\[
 \phi=\left(\frac{43}{3},16,-\frac{25}{2},0,
             \frac{52}{3},-\frac52,-\frac{23}{3}\right).
\]
Its successive gain levels are
$3,10/3,13/3,15/2$, pinning respectively
$\{2,4\},\{1\},\{5,7\},\{3,6\}$; exact primal--dual pairs certify each
leximin stage.  The core point $u$ has $u(S)=36=\phi(S)+8/3$, so $S$ is
binding.  For the star LP, unit weights on coalitions
$\{1,2,3,6,7\}$ and $A$, unit weight on the floor of player $4$, and grand
coefficient $-2$ give an exact dual upper bound $x(S)\leq100/3=\phi(S)$.
Hence the star opening is zero; the corresponding residual vector on $S$ is
$(1,2,1)$.

For trace extension, the restricted witness
$(y_1,y_3,y_5,y_6,y_7)=(13,-5,14,-8,-5)$ attains $y(P)=-5$; unit dual
weights on $\{3,5,6\}$ and $\{1,5,7\}$ certify optimality.  On the full
face, the witness
\[
 (13,47/3,-38/3,1/3,15,-4/3,-5)
\]
attains $x(P)=-6$.  Half-unit dual weights on
$\{3,5,6\},\{1,3,5,6\},\{2,5,7\},\{1,4,5,7\}$, together with grand and
face coefficients $-1/2$ and $-3/2$, certify that this is the full-face
maximum.  This proves \eqref{eq:lx-star-trace-counterexample}; the exact split
LP gives $\Gamma_A=2$.
\end{proof}

The counterexample was missed by broad but non-adaptive screens.  Across $10{,}395$ arbitrary
codimension-two faces of $99$ nondegenerate exact seven-player lower-envelope
games, $86$ restricted optima did not extend, with exactly certified gaps of
$3$ and $6$ in the two control corpora.  Yet none of those gaps occurs among
the $501$ faces conditioned on star failure.  The largest exact ratio of trace
gap to star opening in the first control is $5$, so the tempting quantitative
bound ``gap at most twice the opening'' was already false.  Six rational affine bridges from a genuine three-level
failure to an exact trace-gap control tell the same story: every zero-opening
grid point has zero gap, while the first positive gap appears only at
$t=17/20$ or $9/10$, after star opening is already positive.  No gap occurs in
the six exact three-player failure cases or in five additional float
multilevel candidates from $1{,}484$ nondegenerate targeted games;
exactification sent the sampled extremals from the latter set to the two-level
branch.  A further $1{,}000$-game local mutation screen retained $54$ genuine
three-level failures with a unique minimum residual.  All $54$ have zero
trace gap and positive face-target margin; two rational extremals certify
target margins $15/4$ and $2$.

The decisive search was instead transverse to the exact positive equality
$\Delta^{\rm NL}_{A,5}=\eta_S(\phi)=1$.  Of $252$ two-scale, sum-preserving
generator deformations, all retained a returned unit-price cover and positive
normal gap, while $88$ violated \eqref{eq:lx-star-to-normal} in floating
arithmetic.  The ten closest records all recomputed to exact violations.  On
the rational path in Proposition~\ref{prop:lx-star-trace-counterexample}, the
normal gap stays $1$ and $\Gamma_A$ stays $2$; the exact star opening is
$1/40$ at $\varepsilon=13/20$ and reaches zero at $\varepsilon=2/3$.
Therefore star failure alone does not imply trace extension, and this
unqualified trace route is closed.  The counterexample has residuals
$(1,2,1)$, however, so it belongs to the two-level branch already settled by
Proposition~\ref{prop:lx-two-level-capacity}.  It does not refute the
unresolved genuinely multilevel program.  Proposition~\ref{lem:lx-three-player-order}
then suggests the sharp repair.

\newpage
\begin{conjecture}[multilevel star-conditioned trace extension]
\label{conj:lx-multilevel-trace-extension}
Let $v$ be exact, let $|S|=3$, suppose $(\star_S)$ fails at the leximin point,
and suppose no failure cover has at most two residual levels on $S$.  Let $k$
be the unique minimum-residual player supplied by
Proposition~\ref{lem:lx-three-player-order}, put
$A=(N\setminus S)\cup\{k\}$ and $P=A\setminus\{k\}$.  Then equality holds in
\eqref{eq:lx-trace-extension}.
\end{conjecture}

This is now a seven-case exact frontier.  An exact whole-face audit
now verifies the quantifier in the conjecture, rather than merely inspecting
one returned cover: for each of the three residual pairs, adding the equality
$r_i=r_j$ to the rational failure-cover polyhedron is infeasible in all seven
cases.  Hence every cover has three distinct levels.  Their strict residual
orders are fixed, and after normalizing the minimum residual to one, the exact
minimum gaps to the other two coordinates are one and two.  The original six
cases have $\Gamma_A=0,0,0,-3,0,0$, the seventh case in
Proposition~\ref{prop:lx-three-player-face-target-counterexample} has zero
trace gap, and the $54$ genuine three-level local mutations all have zero
trace gap.  Thus the seventh case falsifies the stronger face target while
continuing to satisfy this trace-extension conjecture.  The earlier
positive-trace-gap counterexample passes the control check: one of its pair-tie
systems is feasible, so it is genuinely outside the conjectured branch.

The new counterexample shows that the multilevel hypothesis cannot simply be
deleted.  A targeted
tie-splitting audit makes $378$ exact sum-preserving deformations around the
counterexample.  Of these, $246$ remain binding star failures; every recovered
failure cover still has exactly two residual levels, and none enters the
three-level chamber.  A second transverse screen makes $756$ deformations
around two distinct whole-face multilevel fixtures.  It retains $181$ binding
star failures and $164$ returned three-level unique-minimum cases, with no
positive trace gap; the extremum from each fixture recomputes exactly to gap
zero and remains whole-face genuinely multilevel.  Thus the positive trace gap
appears locally confined to the already-settled boundary.

A separate $3{,}528$-deformation screen targeted the high-coordinate sign
rather than trace gap.  It retained $1{,}275$ binding star failures and
$1{,}146$ returned three-level failures.  Two had negative float
high-coordinate margins; the worst recomputed exactly to the $-2/3$ face and
coordinate margins in
Proposition~\ref{prop:lx-three-player-face-target-counterexample}.  This closes
the stronger face-target route but leaves both the minimum-residual carrier
inequality and multilevel trace extension intact in the new case.

The new criterion identifies the missing step more sharply.  The exact
$\Gamma_A$ values in the six three-level failures are
$0,0,0,-3,0,0$.  By contrast, the exactly certified worst-gap points in the two
controls have $\Gamma_A=3$ and $6$, while the separate ratio-$5$ control has
$\Gamma_A=5$ even though the first returned restricted optimum has excess
$6$.  All have zero two-player-trace excess, as
Proposition~\ref{prop:lx-two-player-trace} requires.  The counterexample shows
that star failure need not force $\Gamma_A\leq0$ without the multilevel
hypothesis in Conjecture~\ref{conj:lx-multilevel-trace-extension}.

A tempting point-free strengthening is false even inside the hard corpus.
Maximize each marginal $v(B\cup\{p\})-v(B)$ separately over traces $B$ that
contain $k$ and no other member of $S$.  In one exact three-level star failure
the two maxima exceed $d_A$ by $3$, although $\Gamma_A=0$: the selected
restricted optimum places exactly compensating slack on one offending trace.
Thus a universal marginal-worth inequality cannot prove trace extension;
selection inside the restricted optimal face, equivalently the dual balance
\eqref{eq:lx-split-obstruction-balance}, is essential.

The compatible-cover certificate gives proof-producing decompositions for all
six exact three-level cases.  Every optimal split-dual face contains a member
with $\alpha-1=\beta c$, and every minimum-cover support has exact
common-tightness gap zero.  Three cases have $c=0$ and hence follow from
singleton exactness directly; the three $c=1$ cases have exact rational common
core witnesses.  In the sixth case the dual first returned by the LP does not
have the coefficient match, but an exact search of its optimal-dual face finds
the compatible representative $\alpha-1=\beta=1$.  No optimal-dual face of a
positive-gap control contains a compatible member.  The worst-gap controls
also have common-tightness gap $2$; the ratio-$5$ control has a common-tight
minimum support but fails coefficient compatibility throughout its optimal
split-dual face.  Primal tangent minimization gives the same separation: the
base parameter minimizes \eqref{eq:lx-minimum-dual-tangent} in all six hard
cases, whereas the three distinct positive-gap controls have exact tangent
improvements $1$, $6$, and $4$, attained at $t=1$, $11$, and $4$.
The stronger demand that the split weights dominate
$\beta\kappa$ coefficientwise is feasible in only one of the six
hard cases, so the signed common-tightness argument is genuinely needed.  The
counterexample has $c=1$, normal-lift gap $1$, no compatible optimal split
dual, and exact tangent improvement $2$.  Hence neither tangent non-descent
nor the minimum-normal lift follows from star failure and the leximin strata
alone.  Conjecture~\ref{conj:lx-multilevel-trace-extension} excludes the
counterexample's two-level residual pattern; proving or breaking that repaired
statement is now the precise face-lifting target.

The cover-independent audit reveals a substantially simpler target.  It uses
the full minimum-residual class at once.

\begin{conjecture}[minimum-residual-class carrier]
\label{conj:lx-class-carrier}
Let $v$ be exact and let $(\star_S)$ fail at $\phi$.  If the failure is not
already covered by Proposition~\ref{prop:lx-two-level-capacity}, then there is
a failure cover with minimum-residual class
$K=\{i\in S:r_i=r_0\}$ and $P=N\setminus S$ such that
\[
 [\phi(P\cup K)-v(P\cup K)]
 \leq[\phi(K)-v(K)],
 \tag{KC}\label{eq:lx-class-carrier}
\]
or equivalently
\[
 v(P\cup K)-v(K)\geq\phi(P).
\]
\end{conjecture}

This single inequality would close the multilevel capacity branch immediately:
apply Proposition~\ref{prop:lx-carrier-capacity} with $C=K$.  It requires no
singleton selector and is insensitive to how payoffs are ranked inside $K$.
A deduplicated exact audit verifies \eqref{eq:lx-class-carrier} for all $669$
stored multilevel game--coalition pairs ($702$ records): $644$ inequalities
are strict and $25$ are equalities, with margins ranging from $-28$ to $0$.
This includes all $505$ adversarial mutation pairs on which the fixed-cover
escape test mostly fails.  Thus \eqref{eq:lx-class-carrier}, rather than a
singleton rule or residual-gap escape, is now the sharp computational theorem
target.  Deterministic replay of the historical random streams recovered and
verified the discarded point generators for all $25$ equality records.  They
represent $23$ games, and every one is the lower envelope of exactly two
\emph{distinct} common-sum points; the nominal three-generator records merely
repeat one point.  Hence every non-strict case in this exact corpus already
lies in the model class settled by Proposition~\ref{prop:lx-two-generator}.
Outside that proved boundary subcase, all audited class certificates are
strict.

Lemma~\ref{lem:lx-generator-compensation} gives an exact-game face reduction;
the finite-model audit is one way to expose its witnesses.  Among the $530$
deduplicated class-carrier records that retain their point generators, $461$
are direct: some generator attaining
$v(P\cup K)$ already has $a(P)\geq\phi(P)$.  Each of the remaining $69$
records has exactly three distinct generators, a unique $P\cup K$ minimizer,
$|K|=3$, and a strictly positive trace-compensated margin between $19/4$ and
$8$.  All $69$ arise from one seeded target family.  These counts are exact
finite evidence, not a three-generator theorem, but they isolate the only
nontrivial generator mechanism in the embedded corpus.

The quantifier over failure covers in Conjecture~\ref{conj:lx-class-carrier}
is essential.  The next exact three-generator family has failure covers in
three residual chambers.  On an open parameter interval, two chambers select
a class that violates \eqref{eq:lx-class-carrier}, while the third selects the
unique carrier class.  Thus no distinguished output of a failure-cover LP can
be used without a separate selection argument.

\begin{proposition}[failure-cover selection is essential]
\label{prop:lx-cover-selection}
For $0\leq\kappa\leq1$, let $v^\kappa$ be the lower envelope of
\[
\begin{aligned}
 p^\kappa&=(-4,3-\kappa,27,\kappa,10,-8),\\
 q^\kappa&=(-4,1-\kappa,29,-2+\kappa,12,-8),\\
 u^\kappa&=(-\kappa,-1,22,\kappa,10,-3),
\end{aligned}
\qquad
v^\kappa(T)=\min\{p^\kappa(T),q^\kappa(T),u^\kappa(T)\}.
\]
Then $v^\kappa$ is exact and
\[
 \phi^\kappa=\left(-\frac52-\frac\kappa2,
 \frac12-\frac\kappa2,\frac{51}2-\frac\kappa2,-1+\kappa,
 11,-\frac{11}2+\frac\kappa2\right).
\]
For $S=\{2,3,4,6\}$ and $P=\{1,5\}$, the coalition $S$ is binding by
$5/2-\kappa/2$, but $(\star_S)$ fails throughout the interval.  Moreover
every normalized failure cover has, for some
\[
 a\geq1,\qquad t-a\geq1,
\]
the residual vector
\[
 (r_2,r_3,r_4,r_6)=(a,t,t-a,a+t),                       \tag{RC}
 \label{eq:lx-cover-selection-cone}
\]
and every vector in this cone is attained by a failure cover.  Its
minimum-residual class is
\[
 K=\begin{cases}
 \{2\},&t>2a,\\
 \{2,4\},&t=2a,\\
 \{4\},&a+1\leq t<2a.
 \end{cases}
\]
Writing
\[
 \Delta_\kappa(K)
 =[\phi^\kappa(P\cup K)-v^\kappa(P\cup K)]
  -[\phi^\kappa(K)-v^\kappa(K)],
\]
the three possible margins are
\[
 \Delta_\kappa(\{2\})=-\frac32+\frac\kappa2,
 \qquad
 \Delta_\kappa(\{4\})=\Delta_\kappa(\{2,4\})
 =\frac12-\frac\kappa2.
\]
Thus, for $0\leq\kappa<1$, the normalized cover with residuals
$(1,2,1,3)$ selects the noncarrier class $\{2,4\}$, whereas the cover with
residuals $(1,3,2,4)$ selects the carrier class $\{2\}$.  At $\kappa=1$
the former class reaches carrier equality while the latter remains strict.
\end{proposition}

\begin{proof}
All three generators have total $28$.  Each is therefore efficient and
dominates $v^\kappa$ on every coalition, hence belongs to the core; a
generator attains $v^\kappa(T)$ for every $T$, proving exactness.  The
singleton floors are
\[
 m^\kappa=(-4,-1,22,-2+\kappa,10,-8).
\]
We verify the leximin point directly.  If every gain $x_i-m_i^\kappa$ is at
least $s$, the core inequality
$x(\{1,2,3,6\})\geq18-\kappa$ and efficiency imply
$x_4+x_5\leq10+\kappa$, while the floors imply
$x_4+x_5\geq8+\kappa+2s$.  Thus $s\leq1$, and equality pins
$(x_4,x_5)=(-1+\kappa,11)$.  On that face,
$x(\{3,4,6\})\geq19+\kappa$ gives
$x_1+x_2\leq-2-\kappa$; hence the next gain is at most
$(3-\kappa)/2$, with equality pinning
$(x_1,x_2)=(-5/2-\kappa/2,1/2-\kappa/2)$.  On the resulting face,
$x(\{1,3,4\})\geq22$ gives $x_6\leq-11/2+\kappa/2$, so the next gain is at
most $5/2+\kappa/2$.  Equality pins $x_6$, and efficiency fixes
$x_3=51/2-\kappa/2$.  Direct affine endpoint substitution shows that the
displayed point is in the core throughout $0\leq\kappa\leq1$, so these
successive bounds prove the formula for $\phi^\kappa$.

The proper coalitions tight at $\phi^\kappa$ are exactly
\[
 134,\quad125,\quad45,\quad1236,\quad346,
\]
where strings denote player sets.  For each other proper coalition, one fixed
generator lies strictly below $\phi^\kappa$ at both parameter endpoints, so
the same affine strict inequality holds throughout the interval.  The exact
verifier records all $62$ selectors.  The multipliers
\[
 \lambda_{134}=1,\qquad \lambda_{125}=2,
 \qquad\lambda_{45}=1,\qquad \mu=3
\]
cover the two players in $P$ exactly $\mu$ times and leave residuals
$(1,2,1,3)$ on $S$.  If $d$ were a direction in $(\star_S)$, then
\[
 r\mathbin{\cdot}d_S
 =\mu d(N)-\sum_T\lambda_Td(T)\leq0,
\]
but $d_i\geq0$ on $S$ and $d(S)>0$ make the left side positive.  Hence
$(\star_S)$ fails.  Also $v^\kappa(P)=6$, so every core point has
$x(S)=28-x(P)\leq22$; the generator $p^\kappa$ attains equality, while
$\phi^\kappa(S)=39/2+\kappa/2$.  This proves the binding margin
$5/2-\kappa/2$.

It remains to describe the whole cover cone.  Put multipliers $A,B,C,D,E$
on the five tight coalitions in the displayed order.  Exact coverage of
players $1$ and $5$ gives
\[
 \mu=A+B+D=B+C,
 \qquad\text{so}\qquad C=A+D.
\]
Setting $a=A$ and $t=B-E$ now gives
\[
 (r_2,r_3,r_4,r_6)=(a,t,t-a,a+t).
\]
The normalization $r_i\geq1$ reduces to $a\geq1$ and $t-a\geq1$.
Conversely, every such pair is realized by taking
$D=E=0$, $A=C=a$, and $B=t$.  Comparing the four coordinates yields the
three stated residual chambers.  Direct substitution in the worth function
gives the displayed affine values of $\Delta_\kappa(K)$.  Finally,
$(a,t)=(1,2)$ and $(1,3)$ give the two explicit covers claimed.
\end{proof}

A dedicated screen of $500$ six-player, three-generator nearby instances
retained $496$ games with three distinct generators.  It found $30$
multilevel binding star failures and one positive margin for the cover
returned by the search LP; exactification produced the $\kappa=0$ member of
Proposition~\ref{prop:lx-cover-selection}.  An exact prescribed-class audit
then tested every proper class and found the favorable $K=\{2\}$ cover above.
Two targeted mutation sweeps around that member retained $1{,}998$
nondegenerate games and found $28$ multilevel binding star failures.  The
returned cover failed in $16$, but full chamber enumeration found a carrier
cover in every case; the closest exact mutation is the $\kappa=1$ endpoint.
These screens are finite discovery evidence; the proposition itself is
analytic and exhaustive over the full cover cone for every parameter value.

A separate targeted screen of $3{,}000$ four-generator
lower-envelope games found $1{,}123$ star failures and $25$ genuinely
multilevel failures, with no positive class-carrier margin; the closest
boundary case was recomputed exactly and has margin zero.  The broader screen
is discovery evidence only, whereas the $669$-pair audit is fully exact.
The boundary admits an exact parametric analysis.

\begin{proposition}[class-carrier boundary coupling]
\label{prop:lx-class-boundary}
For $0\leq\varepsilon\leq1/3$, let $v^\varepsilon$ be the lower envelope of
the four common-sum points
\[
\begin{aligned}
 p&=(1,1,1,2,1,2),&
 q^\varepsilon&=(0,1,2-\varepsilon,0,4+\varepsilon,1),\\
 r&=(2,1,2,1,2,0),&
 s&=(1,2,0,2,1,2).
\end{aligned}
\]
Put $S=\{2,4,5,6\}$, $P=\{1,3\}$, and $K=\{4,6\}$.  Then
$v^\varepsilon$ is exact,
\[
 \phi(v^\varepsilon)=(1,2,1,1,2,1),
\]
the class-carrier failure margin is exactly
\[
 [\phi(P\cup K)-v^\varepsilon(P\cup K)]
 -[\phi(K)-v^\varepsilon(K)]=\varepsilon,
\]
and the star opening is exactly
\[
 \eta_S(\phi)=3\varepsilon.
\]
Thus crossing the class-carrier boundary in this exact family restores
$(\star_S)$ at three times the crossing rate.
\end{proposition}

\begin{proof}
The generators have common total $8$, so their lower envelope is exact.  The
singleton floors are $(0,1,0,0,1,0)$.  Direct endpoint substitution shows
that $\phi=(1,2,1,1,2,1)$ belongs to the core throughout the interval.  Its
six gains all equal one, while the total gain is six; hence it is the leximin
point.  Moreover,
\[
 v^\varepsilon(K)=1,\qquad
 v^\varepsilon(P\cup K)=3-\varepsilon,\qquad
 \phi(P)=2,
\]
which proves the carrier-margin formula.

The allocation
\[
 y^\varepsilon
 =\phi+\varepsilon(-1,0,-2,1,0,2)
\]
belongs to the core for $0\leq\varepsilon\leq1/3$, dominates $\phi$ on $S$,
and satisfies $y^\varepsilon(S)=6+3\varepsilon$.  Core membership is a finite
affine endpoint check: for each proper coalition the same one of the four
generators lies weakly below $y^\varepsilon$ at both endpoints, hence
throughout the interval.  The standalone exact verifier records all $62$
selectors.

For the reverse inequality, let $x$ be any core point with
$x_i\geq\phi_i$ for $i\in S$.  The three core constraints and two coordinate
floors
\[
 x_{34}\geq2-\varepsilon,\quad
 x_{145}\geq4,\quad
 x_{136}\geq3-\varepsilon,\quad
 x_2\geq2,\quad x_5\geq2
\]
give, after taking multiplicities $(1,1,2,2,1)$,
\[
 3x(P)+2x(S)=24-x(S)\geq18-3\varepsilon.
\]
Thus $x(S)\leq6+3\varepsilon$, and $y^\varepsilon$ attains equality.
\end{proof}

The recovered two-generator boundary has a second exact local test.  Six of
the $25$ equality records had stored a duplicated generator.  A
$1{,}620$-direction single-copy splitting screen found positive
minimum-class margin in exactly those six records.  In the worst screened
direction for each record, exact constraint-generation certificates give
class margin $1/100$ and star opening $1/50$, hence the same ratio $2$ in all
six cases.  These are class-specific comparisons: another carrier has margin
zero in every deformation, so the global $\gamma_S$ defined above remains
zero.  The calculation therefore supplies further boundary coupling for
\eqref{eq:lx-class-carrier}, not an improved coefficient in
\eqref{eq:lx-quantitative-carrier}.  The exact leximin verifier begins with
singleton and co-singleton core constraints and adds every violated omitted
coalition until its rational optimizer satisfies the full core.

Proposition~\ref{prop:lx-two-level-capacity} explains an important part of
this structure exactly.  If the residuals have only the levels $1$ and
$R>1$, the entire minimum-residual class
$C=\{i\in S:r_i=1\}$ is itself the required face certificate; no
crossing-cover construction is needed.  The rare same-stratum pair exceptions
to singleton certification are examples with $|C|=2$.  Accordingly, all
unresolved certificate geometry lies in duals with at least three distinct
residual levels.  Proposition~\ref{prop:lx-carrier-capacity} removes a second
large class: it is enough that an off-side carrier have no more slack than its
trace inside $S$; tightness of the carrier is a particularly simple sufficient
condition.  In particular, the
three-level witness \texttt{WITNESS\_KEY\_FALSE\_n6} is certified by
$T=\{1,2,6\}$ and $C=\{1\}$ (player labels one-indexed), even though the
union-of-cover certificate described above fails.

Proposition~\ref{prop:lx-capacity} removes the frequency language from the
logical frontier. Every pair $(v,S)$ falls into the exact dichotomy
\[
 d^*_v(S)\leq\sigma_S
 \quad\text{or}\quad
 d^*_v(S)>\sigma_S.
\]
The first side admits no binding exact bump and is finished by
Lemma~\ref{lem:lx-nonbinding}. On the second side a binding exact bump exists,
and the missing structural statement is precisely
\[
 d^*_v(S)>\sigma_S\quad\Longrightarrow\quad(\star_S).
\]
Proposition~\ref{prop:lx-two-level-capacity} proves this implication whenever
the failure dual has at most two residual levels, and
Corollary~\ref{prop:lx-pair-capacity} covers every $|S|=2$.  The remaining
capacity frontier therefore requires both $|S|\geq3$ and a genuinely
three-or-more-level dual.  After
Proposition~\ref{prop:lx-carrier-capacity}, a counterexample to the capacity
implication must additionally have no off-side slack-dominating carrier.  The
singleton/pair certificate statistics motivate how to attack those cases, but
neither their observed percentages nor sampling rarity enters the conditional
theorem.

The sharpened computational target is correspondingly narrow.  The broad
slack-domination condition covered all $141$ saved hard star-failure cases at
$n=6,7$ and all $744$ fresh binding star failures in a screen of exact
lower-envelope games at $n=4,5,6$, with no exception.  The new class-carrier
audit finds a successful full minimum-residual class in all $669$ distinct
exact pairs.  Proposition~\ref{prop:lx-cover-selection} shows that this class
need not come from an arbitrarily returned failure cover.  This remains
evidence, not a theorem, but proving
Conjecture~\ref{conj:lx-class-carrier} would close that branch with a single
application of Proposition~\ref{prop:lx-carrier-capacity}.
Proposition~\ref{prop:lx-two-generator} proves the broader existence of some
slack-dominating carrier for every two-generator lower-envelope game.  Thus
within the finite lower-envelope
model class, any counterexample must use at least three genuinely active
generators; this is a structural restriction, not a sampling observation.

The obvious three-generator escape from the sharp two-generator boundary is
itself blocked, but for an instructive reason.  In the exact equality fixture,
let $S=\{1,2,3,4,6\}$ and transfer $a=1/10$ from the off-side player~5 of the
second generator equally among the five members of $S$, retaining both original
generators.  All three generators are genuinely active.  The leximin point
remains $(1,4,3,5,7,4)$ and the minimum off-side slack comparison becomes
$1/50>0$, so the carrier condition has indeed been crossed.  However,
$(\star_S)$ is simultaneously restored: the core point
\[
 (26/25,4,153/50,253/50,171/25,4)
\]
dominates $\phi$ on $S$ and has $S$-sum $429/25>17$.  Thus the star opening is
$4/25$, eight times the carrier escape.  This is an exact boundary calculation,
not a general theorem; a finite simplex-grid deformation screen found the same
coupling throughout the tested local family.  It suggests proving the carrier
conjecture contrapositively: show that escaping every off-side carrier inequality
constructs a core direction nonnegative on $S$, rather than trying to track one
distinguished minimizing generator.

The carrier lift suggests a sharp quantitative form of that contrapositive:
\begin{equation}
 \eta_S(\phi)\geq\gamma_S(\phi).
 \tag{Q}\label{eq:lx-quantitative-carrier}
\end{equation}
Proposition~\ref{prop:lx-quantitative-n7-counterexample} disproves (Q) in
unrestricted dimension, while Proposition~\ref{prop:lx-quantitative-through-six}
proves it through six players.  In view of
Proposition~\ref{prop:lx-exact-core-closure}, this sharp finite-dimensional
theorem applies to every $\psi\in\core(v)$, not only to $\psi=\phi$.
Proposition~\ref{prop:lx-quantitative-pair} proves it for $|S|=2$, and the
coefficient one is already sharp there.  Proposition~\ref{prop:lx-codimension-one-uniform-lift}
proves the stronger uniform lift when $|N\setminus S|=1$; hence
Corollary~\ref{cor:lx-quantitative-through-four} settles (Q) for every exact
game with at most four players, and
Proposition~\ref{prop:lx-quantitative-through-five} closes the remaining
three-versus-two case by exact dual enumeration, even under the weaker
assumption of total balancedness.  Finally,
Proposition~\ref{prop:lx-quantitative-through-six} closes both new
six-player splits by complete minimal-balanced enumeration.  The
three-generator lower envelope in
Proposition~\ref{prop:lx-quantitative-n7-counterexample} then violates (Q) at
seven players at its leximin point with $(\gamma,\eta)=(3/4,0)$ (and at the
core reference zero with $(1,0)$), making the six-player theorem sharp.
Earlier, six exact lower-envelope screens
covering $3{,}650$ games and $35{,}629$ binding pairs $(v,S)$, all $746$ cases
with $\gamma_S(\phi)>0$ satisfied (Q); the minimum in every screen was
recomputed exactly and equalled one.  These counts are now historical
cross-checks, not part of the proof.  The
exactness hypothesis is load-bearing even for the coefficient: a balanced
non-exact five-player game has leximin point
$(-3/2,10,6,17/2,-10)$ and, at $S=\{2,3,4\}$,
$\gamma_S=7/2$ but $\eta_S=1/2$.  Thus the dimension-free target is a
convex-combination statement for the carrier-tight points supplied by
Lemma~\ref{lem:lx-carrier-lift}, not an inequality valid on arbitrary
balanced cores.

Earlier independent screens anticipated the theorem and tested whether it
depended on the leximin reference.  At the
leximin point, $10{,}000$ four-player games yielded $4{,}473$ positive-$\gamma$
three-player targets, including exact equality, and $3{,}000$ five-player
games yielded $3{,}362$ positive-$\gamma$ targets of size at least three; none
had $\eta/\gamma<1$, and the exact five-player minimum was $91/83$.  More
strikingly, the leximin hypothesis was then discarded: $10{,}000$ exact
four-player core references yielded $19{,}521$ positive-$\gamma$ cases with
exact minimum one, while $6{,}400$ exact five-player references yielded
$3{,}732$ positive-$\gamma$ targets of size at least three, again with no
violation.  Finally, $800$ exact references drawn from four saved five-player
games with genuinely non-large cores produced $729$ positive-$\gamma$ targets
of size at least three, no violation, and exact minimum ratio one.  Thus the
pattern is not explained merely by the implication from exactness to largeness
that holds in smaller player sets.  These finite campaigns do not add logical
scope to Propositions~\ref{prop:lx-quantitative-through-five} and
\ref{prop:lx-quantitative-through-six}, but they helped identify the
selection-free exact-core closure inequality
\eqref{eq:lx-convex-gamma}--\eqref{eq:lx-convex-eta}.

The common-coordinate strengthening
$\kappa_S(\psi)\geq\gamma_S(\psi)/|S|$ also survived dense, sparse,
generator, and non-large exact screens in the first open regime
$|S|=3$, $|N\setminus S|=2$.  The largest generator campaign contained
$60{,}418$ positive-$\gamma$ targets and had exact minimum normalized ratio
$47/45$; the non-large audit attained equality.  By contrast, a targeted
$26{,}471$-case screen drove the corresponding outside-slack ratio down to
$3/112$ while retaining carrier-normalized ratio $3/2$.  Thus a proof cannot
replace the carrier slack system by a single outside-coalition cap; the exact
dual in Proposition~\ref{prop:lx-quantitative-orthant-dual} is the sharper
remaining object.

At six players the totally balanced analogue already fails, with
$(\gamma,\eta)=(1,0)$, but exact closure removes the gap.  The resulting
three-block identity, Lemma~\ref{lem:lx-three-block-exact-closure}, first
extends to the complete unit-triple layer of
Proposition~\ref{prop:lx-n6-unit-triple-layer}.  More decisively,
Lemma~\ref{lem:lx-eta-minimal-balanced} converts every relevant eta-dual
vertex into a minimal balanced collection containing $S$.  Exhausting all
$1{,}489$ symmetry orbits and certifying every exact-cone optimum proves the
quantitative carrier inequality for every exact game through six players
(Proposition~\ref{prop:lx-quantitative-through-six}).  The earlier direct
sample of $505$ orbits is recovered in full by the census and now serves as
an independent cross-check rather than merely finite evidence.  The
three-generator exact game in
Proposition~\ref{prop:lx-quantitative-n7-counterexample} has
$(\gamma,\eta)=(3/4,0)$ at its leximin point, so this finite theorem is
dimension-sharp.

\begin{remark}
Lemma~\ref{lem:lx-C} is a leximin parametric-monotonicity statement; its cousin the
nucleolus is \emph{not} coalitionally monotone, so exactness (Lemma
\ref{lem:lx-interval}) is doing real work and the ``self-defeating within-$S$
transfer'' heuristic is not by itself a proof. For proper coalitions of size at
least three we leave (D), (E$'$) open and record
the conditional theorem: \emph{on exact games, $\phi$ is monotone under proper
nonsingleton bumps provided Lemma~\ref{lem:lx-C} holds}; singleton and
grand-coalition bumps remain outside this reduction.  The singleton case is
settled positively by Theorem~\ref{thm:lx-singleton}, while the grand case is
settled negatively by Theorem~\ref{thm:lx-grand-fails}; the necessity direction
(Theorem~\ref{thm:lx-necessity}) is also unconditional.
\end{remark}
